% !Mode:: "TeX:UTF-8"
% !TeX encoding = UTF8
% 用XeLaTeX编译，直接支持中文
%%%%%%%%%%%%%%%%%%%%%%%%%%%%%%%%%%%%%%%%%%%%%%%%%
%% 桂林理工大学本科生学位论文 LaTeX 模版
%% 2020-10-10 v1.0
%% 作者：胡光辉
%% E-mail: guanghui.hu@me.com
%%%%%%%%%%%%%%%%%%%%%%%%%%%%%%%%%%%%%%%%%%%%%%%%%

\documentclass{GLUTthesis}
\usepackage{lscape}
\usepackage{float}
\usepackage{multicol}
\addbibresource{GLUTthesis.bib}

%!TEX root = ../GLUTthesis.tex
% 文章信息
\titlestyle{论文}            % 设计/论文/报告，三选一
\titlecn{非线性 Schr\"odinger 方程数值模拟}   % 论文题目，如果过长，分两行写，第一行写这里
\titlecntwo{单光孤子的周期性演化}  % 题目过长，分两行后的第二行写这里
\titleen{Numerical simulation of nolinear Schr\"odinger equation\\Periodic evolution of single optical soliton}  % 英文题目
\major{应用物理学}             % 专业
\class{2018-2班}              % 班级
\author{韩旭东}                 % 姓名
\authoren{HAN Xu-dong}          % 姓名拼音，注意大小写
\supervisor{王恒}            % 你的导师
\supervisoren{WANG Heng}   % 你导师姓名的拼音，注意大小写
\subsupervisor{}             % 有第二导师的话，没有留空
\department{理学院}           % 学院
\studentid{3182090711240}      % 你的学号
\donedata{\zhtoday}      % 论文完成日期

\begin{document}
%%%%%%%%%%%%%%%%%%%%%%%%%%%%%%%%%%%%%%%%%%%%%%%%%%
% 封面
% -----------------------------------------------%
\makecoverpage
\newpage
%%%%%%%%%%%%%%%%%%%%%%%%%%%%%%%%%%%%%%%%%%%%%%%%%%
% 前置部分的页眉页脚设置
% -----------------------------------------------%
% 前置部分用罗马数字单独编连续码（封面除外）。
\pagenumbering{Roman} % 大写罗马字母
\setcounter{page}{1} % 从1开始编号页码

%%%%%%%%%%%%%%%%%%%%%%%%%%%%%%%%%%%%%%%%%%%%%%%%%%
% 声明页
% -----------------------------------------------%
\announcement
\newpage
%%%%%%%%%%%%%%%%%%%%%%%%%%%%%%%%%%%%%%%%%%%%%%%%%%
% 中文摘要
% -----------------------------------------------%
\addcontentsline{toc}{section}{摘\quad 要}
%!TEX root = ../GLUTthesis.tex
% 设置中文摘要
\keywordscn{非线性 Schr\"odinger 方程；光孤子；伪谱方法；分步 Fourier 方法；有限差分法}

\begin{abstractcn}

    非线性 Schr\"odinger 方程不仅是非线性光纤光学的重要问题，也是非线性偏微分方程的重要课题之一。数学上可应用逆散射方法得到非线性 Schr\"odinger 方程的解析解，但是需要求解很复杂的积分方程，因此数值方法求解非线性 Schr\"odinger 方程是有必要的。非线性 Schr\"odinger 方程有一类特殊的解，在传输过程中，光脉冲会经历窄化、分裂、再恢复的周期性演化过程。它就是光孤子，它是色散和非线性效应共同作用的结果。本文用伪谱方法、分步 Fourier 方法和有限差分法实现了非线性 Schr\"odinger 方程的数值求解，得到了单光孤子周期性演化的物理图像。

\end{abstractcn}
\newpage

%%%%%%%%%%%%%%%%%%%%%%%%%%%%%%%%%%%%%%%%%%%%%%%%%%
% 英文摘要
% -----------------------------------------------%
\addcontentsline{toc}{section}{Abstract}
%!TEX root = ../GLUTthesis.tex
\keywordsen{NLSE; optical soliton; pseudospectral method; SSFM; finite difference method}
\begin{abstracten}
    Nonlinear Schr\"odinger equation(NLSE) is not only an important problem in nonlinear fiber optics, but also one of the important topics in nonlinear partial differential equations. The Inverse Scattering Transform Method(IST) can be used to get the analytical solution of nonlinear Schr\"odinger equation In mathematics, but it needs to solve a very complex integral equation. So it is necessary to solve nonlinear Schr\"odinger equation numerically. There is a special kind of solution in nonlinear Schr\"odinger equation. During the transmission process, the optical pulse will undergo a periodic evolution process of narrowing, splitting, and recovery. It is an optical soliton, which is the result of the combined action of dispersion and nonlinear effects. In this paper, pseudospectral method, split-step Fourier method(SSFM) and finite difference method are used to solve the nonlinear Schr\"odinger equation numerically, and the physical picture of the periodic evolution of single optical soliton is obtained.
\end{abstracten}
\newpage

%%%%%%%%%%%%%%%%%%%%%%%%%%%%%%%%%%%%%%%%%%%%%%%%%%
% 目录
% -----------------------------------------------%
{
\renewcommand{\contentsname}{\hfill \bfseries \zihao{3} 目\quad 录\hfill}  
	\renewcommand*{\baselinestretch}{1.2}   % 行间距
    \tableofcontents
}
\newpage
% 去掉页眉章节序号后面的“.”
\renewcommand{\sectionmark}[1]{\markright{\thesection~ #1}}


\renewcommand{\headrulewidth}{1pt}

%%%%%%%%%%%%%%%%%%%%%%%%%%%%%%%%%%%%%%%%%%%%%%%%%%
% 正文和后置部分的页眉页脚设置
% -----------------------------------------------%
% 正文和后置部分用阿拉伯数字编连续码。
\setcounter{page}{1} % 重置页码编号
\pagenumbering{arabic} % 设置页码编号为阿拉伯数字
% 定义页眉页脚 %
\newcommand{\makeheadrule}{%
    \makebox[-3pt][l]{\rule[.7\baselineskip]{\headwidth}{0.5pt}}
    \rule[0.85\baselineskip]{\headwidth}{1.5pt}\vskip-.8\baselineskip}
\makeatletter
\renewcommand{\headrule}{%
    {\if@fancyplain\let\headrulewidth\plainheadrulewidth\fi
     \makeheadrule}}
% -----画单隔线 ------------
\renewcommand{\headrulewidth}{0.5pt}    % 在页眉下画一个0.5pt宽的分隔线
\renewcommand{\footrulewidth}{0pt}      % 在页脚不画分隔线。
% -----设置页眉样式 ---------
\newcommand{\headstyle}{
	\fancyhead[C]{\color{black} \bfseries \texttt{\zihao{-3} 桂\hspace{0.2em}林\hspace{0.2em}理\hspace{0.2em}工\hspace{0.2em}大\hspace{0.2em}学\hspace{0.2em}本\hspace{0.2em}科\hspace{0.2em}毕\hspace{0.2em}业\hspace{0.2em}设\hspace{0.2em}计\hspace{0.2em}{\zihao{2}$\cdot$}\hspace{0.2em}论\hspace{0.2em}文}}
}
% -----设置页脚样式 ---------
\newcommand{\footstyle}{\fancyfoot[C]{\thepage}}
\pagestyle{empty}
\pagestyle{fancy}
\fancyhf{} %清空原有样式
\headstyle
\footstyle
% ------ 定义一种新的格式叫做main ------------
\fancypagestyle{main}{%
    \pagestyle{fancyplain} 
	\fancyhf{} %清空原有样式
	\headstyle
	\footstyle
}

% 可以使用include命令导入tex文件，从而避免过多修改本文件。

% 论文正文是主体，主体部分应从另页右页开始，每一章应另起页。一般由序号标题、文字叙述、图、表格和公式等五个部分构成。

% 重新设置正文行间距，因为前置部分设置时候行间距被改过
\renewcommand*{\baselinestretch}{1.0}   % 几倍行间距
\setlength{\baselineskip}{20pt}         % 基准行间距

% 正文
% %!TEX root = ../GLUTthesis.tex

%为了便于查找和修改，还可以拆分文件为子章节，并通过input导入。


\section{模板结构}

模板文件的结构，如下表所示：
\begin{table}[ht]\centering
\begin{tabular}{r|l|l}
	\hline
	\multicolumn{2}{l|}{GLUTthesis.tex }  & 主文档，在其中填写正文  \\ \hline
content 文件夹 & info.tex  & 论文信息  \\ \cline{2-3} \hline
content 文件夹 & abstractcn.tex & 中文摘要、关键词  \\ \hline
content 文件夹 & abstracten.tex & 英文摘要、关键词  \\ \hline
content 文件夹 & content.tex & 正文  \\ \hline
content 文件夹 &  thanks.tex & 致谢  \\ \hline
content 文件夹 & additional.tex & 附录  \\ \hline
	\multicolumn{2}{l|}{images 文件夹}                  & 存放图片文件.                   \\ \hline
	\multicolumn{2}{l|}{GLUTthesis.cls}             & 定义文档格式的 class file，不可删除 \\ \hline
	\multicolumn{2}{l|}{GLUTthesis.bib}             & 参考文献存放 \\ \hline
	\multicolumn{2}{l|}{gbt7714.sty}             & 参考文献格式，不可删除 \\ \hline
	\multicolumn{2}{l|}{gbt7714-unsrt.bst}             & 参考文献格式，不可删除 \\ \hline
\end{tabular}
\end{table}

无需也不要改变、移动上述文档的位置。

 \subsection{使用步骤}

1、进入 content 文件夹,  打开 info.tex 填写论文相关信息；
打开 abstractcn.tex、abstracten.tex 这两个文档，分别填写中、英文摘要、关键词；打开 thanks.tex 填写致谢；打开 content.tex 填写正文部分；打开 addtional.tex 填写附录内容。

2、将图片放入 images 文件夹。

3、将参考文献信息录入 GLUTthesis.bib，并在正文中正确引用。

4、使用 XeLaTeX 编译。具体见 \ref{sec-compile} 节。



\subsection{编译方法} \label{sec-compile}

使用 XeLaTeX 编译，直接生成~pdf 文件。


\subsection{其他}
\subsubsection{其他1}
无。
\subsubsection{其他2}
无。
\newpage 


%%%%%%%%%%%%%%%%%%%%%%%%%%%%%%%%%%%%%%%%%%%%%%%%%%
\section{字体操作}
\subsection{字体调节}

\begin{tabular}{ll}
	\verb|\songti|   & {\songti 宋体}   \\
	\verb|\bfseries|    & {\bfseries 粗宋体}    \\
	\verb|\sffamily| & {\sffamily 黑体} \\
	\verb|\bfseries\sffamily| & {\bfseries\sffamily 粗黑体} \\
  \verb|\ttfamily|   & {\ttfamily 楷体}   \\
  \verb|\bfseries\ttfamily|   & {\bfseries\ttfamily 粗楷体}   \\
  \verb|\itshape|   & {\itshape 斜宋体}   \\
  \verb|\bfseries\itshape|   & {\bfseries\itshape 粗斜宋体}   \\

\end{tabular}
\textbf{}

\subsection{字号调节}
字号命令: \verb|\zihao| \index{zihao}

\begin{tabular}{ll}
\verb|\zihao{0}| &\zihao{0}  初号字 English \\
\verb|\zihao{-0}|&\zihao{-0} 小初号 English \\
\verb|\zihao{1} |&\zihao{1}  一号字 English \\
\verb|\zihao{-1}|&\zihao{-1} 小一号 English \\
\verb|\zihao{2} |&\zihao{2}  二号字 English \\
\verb|\zihao{-2}|&\zihao{-2} 小二号 English \\
\verb|\zihao{3} |&\zihao{3}  三号字 English \\
\verb|\zihao{-3}|&\zihao{-3} 小三号 English \\
\verb|\zihao{4} |&\zihao{4}  四号字 English \\
\verb|\zihao{-4}|&\zihao{-4} 小四号 English \\
\verb|\zihao{5} |&\zihao{5}  五号字 English \\
\verb|\zihao{-5}|&\zihao{-5} 小五号 English \\
\verb|\zihao{6} |&\zihao{6}  六号字 English \\
\verb|\zihao{-6}|&\zihao{-6} 小六号 English \\
\verb|\zihao{7} |&\zihao{7}  七号字 English \\
\verb|\zihao{8} |&\zihao{8}  八号字 English \\
\end{tabular}

\newpage
%%%%%%%%%%%%%%%%%%%%%%%%%%%%%%%%%%%%%%%%%%%%%%%%%%%

\section{列表环境}

手工编号

（1）图片插入布局，如第\ref{sec.figure}章所示。

（2）XXXXXXXXXX

（3）XXXXXXXXXX

（4）XXXXXXXXXX

罗马编号
\begin{enumerate}[label=(\roman*)]
 \item XXXXXXXXXX
 \item XXXXXXXXXX
 \item XXXXXXXXXX
\end{enumerate}

括号编号
\begin{enumerate}[label=(\arabic*)]
 \item XXXXXXXXXX
 \item XXXXXXXXXX
 \item XXXXXXXXXX
\end{enumerate}

半括号编号
\begin{enumerate}[label=\arabic*)]
 \item XXXXXXXXXX
 \item XXXXXXXXXX
 \item XXXXXXXXXX
\end{enumerate}

小字母编号
\begin{enumerate}[label=\alph*)]
 \item XXXXXXXXXX
 \item XXXXXXXXXX
 \item XXXXXXXXXX
\end{enumerate}

\newpage

%%%%%%%%%%%%%%%%%%%%%%%%%%%%%%%%%%%%%%%%%%%%%%%%%%%%%%%
\section{公式插入}

公式插入示例如公式(\ref{E.example})所示。

薛定谔方程在一维量子问题中的应用：

（1）一维无限深势阱的势能分布

\begin{equation}
\displaystyle U(x)=\left\{
\begin{array}{l@{}}
0\;\;\quad(0<x<a) \\
U_0\quad (x\leq0 \;\mbox{或} \;x\geq a)\\ 
\end{array}\right.
\label{E.example}
\end{equation}

其本征能量与本征态为
\begin{equation}
\displaystyle U(x)=\left\{
\begin{array}{l@{}}
\displaystyle E_n=\frac{\hbar^2k^2}{2m}=\frac{\pi^2\hbar^2}{2ma^2}n^2 \;\;(n=1\mbox{，}2\mbox{，}3\cdots) \\
\displaystyle \psi _n(x)=\sqrt{\frac{2}{a}}\mathrm{sin}\frac{n\pi}{a}x \;\;\;(0\leq x\leq a) \\ 
\displaystyle \psi _n(x)=0\quad\quad\quad\quad\quad(x\leq0 \;\mbox{或} \;x\geq a) \\ 
\end{array}\right.
\end{equation}

（2）一维谐振子的势能函数$\displaystyle U(x)=\frac{1}{2}kx^2=\frac{1}{2}m\omega^2x^2$，谐振子的定态薛定谔方程为

\begin{equation}
\displaystyle -\frac{\hbar^2}{2m}\frac{\partial^2\psi}{\partial x^2}+\frac{1}{2}m\omega^2x^2\psi=E \psi
\end{equation}

其本征能量与本征态为
\begin{equation}
\displaystyle E_n=(n+\frac{1}{2})\hbar \omega \;\;(n=0\mbox{，}1\mbox{，}2\cdots)
\end{equation}
 
\begin{equation}
	\displaystyle \psi _n(\alpha x)=(\sqrt{\pi}2^nn!)^{-\frac{1}{2}}\mathrm{e}^{-\frac{\alpha^2x^2}{2}} H_n(\alpha x) \;\;(n=0\mbox{，}1\mbox{，}2\cdots)
\end{equation}

（3）一维方势垒和隧道效应：方势垒的势能为
\begin{equation}
	\displaystyle U(x)=\left\{
\begin{array}{l@{}}
U_0\quad(0<x<a) \\
0\;\;\quad( x\leq0 \;\mbox{或} \;x\geq a)\\ 
\end{array}\right.
\end{equation}

投射系数为$\displaystyle D=D_0\mathrm{e}^{-2ka}$，其中$\displaystyle k=\sqrt{\frac{2m(U_0-E)}{\hbar^2}}$。

我们还可以轻松打出一个矩阵
\begin{equation}
\bm{A}=\begin{bmatrix} %\bm为数学粗体 \hm为更粗的加重体数学负号
1&2&3&4\\
11&22&33&44\\
\end{bmatrix}
\times\begin{bmatrix}
22&24\\
32&34\\
42&44\\
52&54\\
\end{bmatrix}
\end{equation}

\newpage

%%%%%%%%%%%%%%%%%%%%%%%%%%%%%%%%%%%%%%%
\section{化学方程式} 

化学方程式可以直接采用数学式输入，例如: 
三硝基甲苯(TNT)~C$_6$H$_2$CH$_3$(NO$_2$)$_3$
为白色或苋色淡黄色针状结晶，无臭，有吸湿性。

很明显，将化学方程式作为数学公式输入很复杂，且十分笨拙。所以本模版引入了mhchem宏包将问题简化。~\verb|\ce|~命令用来输入化学方程式。 如：醋中主要是 \ce{H2O}，含有 \ce{CH3C00-}。\ce{^{277}_{90}Th}元素具有强放射性。

 化学反应式例子如下：
 \begin{gather} 
   \ce{2H2 + O2 ->[\text{燃烧}] 2H2O} \\
   \ce{N2 + 3H2 <=>T[高温、加压][催化剂] 2NH3}
 \end{gather}

有机化学式的书写，先简单介绍一下~chemfig~方向的定义，如图~\ref{fig:direction}：

\begin{figure}[ht]
  \centering
    \includegraphics[width=0.2\textwidth]{direction.png}
    \caption{chemfig方向的定义}
    \label{fig:direction}
\end{figure}

举个乙烷的例子：

\begin{center}
  \chemfig{C(-[2]H)(-[4]H)(-[6]H)-C(-[2]H)(-[6]H)-H}
\end{center}

上述乙烷代码如下：
{\verb|\chemfig{C(-[2]H)(-[4]H)(-[6]H)-C(-[2]H)(-[6]H)-H}|}
代码中 (-[X]Y) 内的数字 X 就表示了内容 Y 的位置，其中中括号（[ ]）的位置通常紧跟在化学键（-、= 等）之后。例如2表示的就是向上的方向，4 表示向左，6 表示向右。

再画一个苯环:
\begin{center}
  \chemfig{*6(-=-=-=)}
\end{center}

%%%%%%%%%%%%%%%%%%%%%%%%%%%%%%%%%%%%%%%%%%%%%%%%%%%%%%

\section{图像布局}
\label{sec.figure}
\subsection{单图布局}

    这是一段随机插入的文本，用来填充模板布局，感受模板视觉效果。这是一段随机插入的文本，用来填充模板布局，感受模板视觉效果。这是一段随机插入的文本，用来填充模板布局，感受模板视觉效果。这是一段随机插入的文本，用来填充模板布局，感受模板视觉效果。这是一段随机插入的文本，用来填充模板布局，感受模板视觉效果。这是一段随机插入的文本，用来填充模板布局，感受模板视觉效果。
    
    这是一段随机插入的文本，用来填充模板布局，感受模板视觉效果。这是一段随机插入的文本，用来填充模板布局，感受模板视觉效果。这是一段随机插入的文本，用来填充模板布局，感受模板视觉效果。这是一段随机插入的文本，用来填充模板布局，感受模板视觉效果。这是一段随机插入的文本，用来填充模板布局，感受模板视觉效果。这是一段随机插入的文本，用来填充模板布局，感受模板视觉效果。

单图布局如图\ref{F.glut_single}所示。

\begin{figure}[hbt]
\centering
\includegraphics[width=1.0\textwidth]{banner.jpeg}
\caption{单图布局示例}
\label{F.glut_single}
\end{figure}

\subsection{横排布局}

横排布局如图\ref{F.glut_row}所示。

\begin{figure}[!htb]
    \centering
    \begin{subfigure}[t]{0.24\linewidth}
        \begin{minipage}[b]{1\linewidth}
        \includegraphics[width=1\linewidth]{glut.png}
        \caption{第一张}
        \end{minipage}
    \end{subfigure}
    \begin{subfigure}[t]{0.24\linewidth}
        \begin{minipage}[b]{1\linewidth}
        \includegraphics[width=1\linewidth]{glut.png}
        \caption{第二张}
        \end{minipage}
    \end{subfigure}
    \caption{横排布局示例}
    \label{F.glut_row}
\end{figure}

    这是一段随机插入的文本，用来填充模板布局，感受模板视觉效果。这是一段随机插入的文本，用来填充模板布局，感受模板视觉效果。这是一段随机插入的文本，用来填充模板布局，感受模板视觉效果。这是一段随机插入的文本，用来填充模板布局，感受模板视觉效果。这是一段随机插入的文本，用来填充模板布局，感受模板视觉效果。这是一段随机插入的文本，用来填充模板布局，感受模板视觉效果。
    
    这是一段随机插入的文本，用来填充模板布局，感受模板视觉效果。这是一段随机插入的文本，用来填充模板布局，感受模板视觉效果。这是一段随机插入的文本，用来填充模板布局，感受模板视觉效果。这是一段随机插入的文本，用来填充模板布局，感受模板视觉效果。这是一段随机插入的文本，用来填充模板布局，感受模板视觉效果。这是一段随机插入的文本，用来填充模板布局，感受模板视觉效果。

\subsection{竖排布局}
竖排布局如图\ref{F.glut_col}所示。

\begin{figure}[!htb]
    \centering
    \begin{subfigure}[t]{0.15\linewidth}
        \captionsetup{justification=centering} %ugly hacks
        \begin{minipage}[b]{1\linewidth}
        \includegraphics[width=1\linewidth]{glut.png}
        \caption{第一张}
        \end{minipage}
    \end{subfigure}\\
    \begin{subfigure}[t]{0.15\linewidth}
        \captionsetup{justification=centering} %ugly hacks
        \begin{minipage}[b]{1\linewidth}
        \includegraphics[width=1\linewidth]{glut.png}
        \caption{第二张}
        \end{minipage}
    \end{subfigure}
    \caption{竖排布局示例}
    \label{F.glut_col}
\end{figure}

    这是一段随机插入的文本，用来填充模板布局，感受模板视觉效果。这是一段随机插入的文本，用来填充模板布局，感受模板视觉效果。这是一段随机插入的文本，用来填充模板布局，感受模板视觉效果。这是一段随机插入的文本，用来填充模板布局，感受模板视觉效果。这是一段随机插入的文本，用来填充模板布局，感受模板视觉效果。这是一段随机插入的文本，用来填充模板布局，感受模板视觉效果。
    
    这是一段随机插入的文本，用来填充模板布局，感受模板视觉效果。这是一段随机插入的文本，用来填充模板布局，感受模板视觉效果。这是一段随机插入的文本，用来填充模板布局，感受模板视觉效果。这是一段随机插入的文本，用来填充模板布局，感受模板视觉效果。这是一段随机插入的文本，用来填充模板布局，感受模板视觉效果。这是一段随机插入的文本，用来填充模板布局，感受模板视觉效果。

\subsection{竖排多图横排布局}

\begin{figure}[!htb]
    \centering
    \begin{subfigure}[t]{0.13\linewidth}
        \captionsetup{justification=centering} 
        \begin{minipage}[b]{1\linewidth}
        \includegraphics[width=1\linewidth]{glut.png} \vspace{-1ex} \vfill
        \includegraphics[width=1\linewidth]{glut.png}
        \caption{第一张}
        \end{minipage}
    \end{subfigure}
    \begin{subfigure}[t]{0.13\linewidth}
        \captionsetup{justification=centering} 
        \begin{minipage}[b]{1\linewidth}
        \includegraphics[width=1\linewidth]{glut.png} \vspace{-1ex} \vfill
        \includegraphics[width=1\linewidth]{glut.png}
        \caption{第二张}
        \end{minipage}
    \end{subfigure}
    \caption{竖排多图横排布局}
    \label{F.glut_col_row}
\end{figure}

竖排多图横排布局如图\ref{F.glut_col_row}所示。注意看(a)、(b)编号与图关系。


\subsection{横排多图竖排布局}

    这是一段随机插入的文本，用来填充模板布局，感受模板视觉效果。这是一段随机插入的文本，用来填充模板布局，感受模板视觉效果。这是一段随机插入的文本，用来填充模板布局，感受模板视觉效果。这是一段随机插入的文本，用来填充模板布局，感受模板视觉效果。这是一段随机插入的文本，用来填充模板布局，感受模板视觉效果。这是一段随机插入的文本，用来填充模板布局，感受模板视觉效果。

\begin{figure}[!htb]
    \centering
    \begin{subfigure}[t]{0.3\linewidth}
        \captionsetup{justification=centering} 
        \begin{minipage}[b]{1\linewidth}
        \includegraphics[width=0.45\linewidth]{glut.png}
        \includegraphics[width=0.45\linewidth]{glut.png}
        \caption{}
        \end{minipage}
    \end{subfigure}\\
    \begin{subfigure}[t]{0.3\linewidth}
        \captionsetup{justification=centering} 
        \begin{minipage}[b]{1\linewidth}
        \includegraphics[width=0.45\linewidth]{glut.png}
        \includegraphics[width=0.45\linewidth]{glut.png}
        \caption{}
        \end{minipage}
    \end{subfigure}
    \caption{横排多图竖排布局}
    \label{F.glut_row_col}
\end{figure}

横排多图竖排布局如图\ref{F.glut_row_col}所示。注意看(a)、(b)编号与图关系。

\newpage
%%%%%%%%%%%%%%%%%%%%%%%%%%%%%%%%%%%%%%%%%%%%%%%%%%%%%%%

\section{表格插入示例}
表格如表\ref{T.example}所示，latex表格技巧很多，这里不再详细介绍。

\begin{table}[htb]
  \centering
  \caption{按照学术论文的一般规范，表格为三线表。}
  \label{T.example}
  \begin{tabular}{ccccc}
\toprule
A   &  B  &  C  & D  & E \\
\midrule
1 	& 一 & 壹 & one 	 & 啊\\
2 	& 二 & 贰 & two 	 & 哈\\
3 	& 三 & 叁 & three & 嗯\\
4 	& 四 & 肆 & four	 & 嘿\\
\bottomrule
\end{tabular}
\end{table}

\newpage

%%%%%%%%%%%%%%%%%%%%%%%%%%%%%%%%%%%%%%%%%%%%%%%%%%%%%%
\section{参考文献插入示例}

LaTeX\cite{lamport1994latex}插入参考文献最方便的方式是使用bibliography\cite{hu2020}，大多数出版商的论文页面\cite{lamport1994latex}都会有导出bib格式参考文献的链接，把每个文献的bib放入``GLUTthesis.bib''，然后用bibkey即可插入参考文献。

\newpage
%%%%%%%%%%%%%%%%%%%%%%%%%%%%%%%%%%%%%%%%%%%%%%%%%%%%%%
\section{总结与展望}

XX的XX都存在XX，所以我们XX，本章总结XX。

\subsection{本文工作总结}
在总结和分析已有XX的理论基础上，本文对XX进行了XX，主要工作如下：

%手工编号
（1）图片插入布局，如第\ref{sec.figure}章所示。

（2）XXXXXXXXXX

（3）XXXXXXXXXX

（4）XXXXXXXXXX


\subsection{工作展望}
本课题针对XX，鉴于XXX，对XX进行了提高，但是XXX，所以有如下XX：

（1）目前XX虽然XX，但是XX仍然XX，所以XX仍然是一个值得XX的问题。

（2）随着XX，XX具有XX的问题，仍值得进一步XX。

（3）本课题在XX有了XX，但是XX的XX还存在XX，所以XX。

\newpage


\section{引言}
\subsection{研究背景}
激光器和超低损耗光纤的问世不仅引起了光纤通信的革命，而且拉开了非线性光学的序幕。非线性光学是现代光学很重要的分支，它主要研究强光（主要是激光）与物质相互作用产生的各种非线性效应，比如 Kerr 效应、受激散射效应、多波混频效应和光孤子效应等。1973 年，Hasegawa 和 Tappert 在研究脉冲在光纤中传输的课题时认为，反常群速度色散和非线性效应的相互作用，可以使孤子脉冲在低损光纤中传输，且其传输过程由非线性 Schr\"odinger 方程描述。他们还指出，若在通信系统用光孤子作为载体，可使通信容量比线性通信系统提高 1$\sim$2 个数量级\cite{Hasegawa}。光孤子通信的广阔前景引起人们的兴趣，对它进行了广泛的研究。

理论上对光孤子的研究就是求解非线性 Schr\"odinger 方程。1971 年，Zakharov 和 Shabat 用逆散射方法对非线性 Schr\"odinger 方程进行求解，得到该方程的解析解。然而，逆散射方法求解非线性 Schr\"odinger 方程需要求解复杂的积分方程，因此数值模拟成为研究非线性 Schr\"odinger 方程（光孤子问题）必要的一种手段。除了求解偏微分方程通用的有限差分法外，Hardin 和 Tappert 在 1972 年提出用 分步 Fourier 方法求解非线性偏微分方程\cite{Tappert}；1978 年，Fornberg 和 Whitham 提出用伪谱方法求解非线性偏微分方程\cite{Fornberg}。非线性 Schr\"odinger 方程的数值求解方法有很多，Taha 和 Ablowitz 在 1983 年对非线性 Schr\"odinger 方程的数值方法进行了总结和比较\cite{Taha}。

\subsection{本文工作}
本文的主要工作是在 MATLAB 上应用有限差分、分步 Fourier 方法和伪谱方法这三种数值方法实现非线性 Schr\"odinger 方程的数值求解。求解目标是得到现单光孤子周期性演化的物理图像。在本文的第\ref{sec:NLSE}章，从光纤介质的本构关系出发出发推导了非线性 Schr\"odinger 方程；在本文的第\ref{sec:optical soliton}章，介绍了孤子理论，着重介绍了但光孤子的传输；在本文的第\ref{sec:numerical}章，阐述了用三种数值方法对光孤子效应进行数值模拟的工作；在本文的最后部分，对三种数值方法进行了总结和对比。

\section{非线性 Schr\"odinger 方程}
\label{sec:NLSE}
\subsection{介电响应}

根据 Maxwell 电磁理论，光在光纤中的传输可以用波动方程描述
\begin{equation}
    \nabla^2\vec{E}-\frac{1}{c^2}\frac{\partial^2\vec{E}}{\partial t^2}-\mu_0\frac{\partial^2\vec{P}}{\partial t^2}=\vec{0}\\
\end{equation}

光纤的主要材料是熔融 \ce{SiO_2} 合成的纯石英玻璃，该电介质对光波的介电响应 $\vec{P}(\omega)$ 与光波频率 $\omega$ 有关。光纤的这一特性就是色散效应，最终导致光脉冲不同的频率分量以不同的速度传输。因此在频域中讨论光波问题会更方便，对波动方程进行 Fourier 变换得到三个分量的 Helmholtz 方程
\begin{equation}
    \nabla^2\widetilde{\vec{E}}(\vec{r},\omega)+\varepsilon_r(\omega)\frac{\omega^2}{c^2}\widetilde{\vec{E}}(\vec{r},\omega)=\vec{0}
\end{equation}

短脉冲（脉宽范围约为 10 fs $\sim$ 10 ns）在光纤中传输时，介质对高强度光场的响应是非线性的，通常情况下极化强度 $\vec{P}$ 与电场 $\vec{E}$ 的关系为
\begin{equation}
    \vec{P}=\varepsilon_0\left(\chi^{(1)}\cdot\vec{E}+\chi^{(2)}:\vec{E}\vec{E}+\chi^{(3)}\vdots\vec{E}\vec{E}\vec{E}+\cdots\right)    
\end{equation}
其中，线性极化率 $\chi^{(1)}$ 对介电响应起主要作用，影响光纤的折射率 $n$ 和衰减系数 $\alpha$；由于 \ce{SiO_2} 为对称分子，光纤中的二阶极化率 $\chi^{(2)}=0$；三阶极化率 $\chi^{(3)}$ 及更高阶的极化率是产生非线性光学现象的主要原因。

因此，光脉冲在光纤中传输时，需要综合考虑色散效应和非线性效应。如此复杂的本构关系，势必会增加求解 Helmholtz 方程的难度，那么不妨做如下假设\cite{Agrawal}
\begin{enumerate}[label=(\arabic*)]
    \item 把非线性介电响应处理为线性介电响应的微扰；
    \item 假定光场的偏振态沿光纤长度方向不变；
    \item 假定光场是准单色的，即中心频率 $\omega_0$ 远小于频谱宽度 $\Delta\omega$，$\Delta\omega/\omega_0\ll1$。
\end{enumerate}
这样一来，讨论光在光纤中的传输问题简化为求解偏振方向的  Helmholtz 方程
\begin{equation}
    \nabla^2\widetilde{E}+\varepsilon_r(\omega)k_0^2\widetilde{E}=0
    \label{eq:Helmholtz}
\end{equation}

相对介电常数 $\varepsilon_r$ 可以表达为\footnote{线偏振光场假设下，仅需考虑极化率张量的某一分量。}
\begin{equation}
    \varepsilon_r(\omega)=1+\widetilde{\chi}_{xx}^{(1)}+\varepsilon_{NL}
\end{equation}
其中，线性极化率 $\widetilde{\chi}_{xx}^{(1)}(\omega)$ 是 $\chi_{xx}^{(1)}(t)$ 作 Fourier 变换的结果，非线性介电常数 $\varepsilon_{NL}$ 被处理为常量
\begin{equation}
    \varepsilon_{NL}=\frac{3}{4}\chi_{xxxx}^{(3)}\left|\vec{E}(\vec{r},t)\right|^2    
\end{equation}
可用介电常量 $\varepsilon_r$ 定义折射率 $\widetilde{n}$ 和吸收系数 $\widetilde{\alpha}$
\begin{equation}
    \varepsilon_r=\left(\widetilde{n}+\frac{i\widetilde{\alpha}}{2k_0}\right)^2
\end{equation}
介电常数 $\varepsilon_r$ 中含有与光场有关的非线性部分 $\varepsilon_{NL}$，折射率 $\widetilde{n}$ 和吸收系数 $\widetilde{\alpha}$ 描述如下
\begin{subequations}
    \begin{equation}
        \widetilde{n}=n+n_2|E|^2        
    \end{equation}
    \begin{equation}
        \widetilde{\alpha}=\alpha+\alpha_2|E|^2        
    \end{equation}
\end{subequations}
其中，线性折射率 $n$ 和线性吸收系数 $\alpha$ 与线性极化率 $\widetilde{\chi}_{xx}^{(1)}$ 有关；非线性折射率 $n_2$ 和非线性吸收系数 $\alpha_2$ 与三阶极化率 $\chi_{xxxx}^{(3)}$ 有关。通常忽略石英光纤的非线性吸收系数 $\alpha_2$。
\subsection{脉冲传输}
\label{sec:Schrodinger}
研究短脉冲在光纤中的传输问题就是求解 Helmholtz 方程(\ref{eq:Helmholtz})，试探解分离变量如下
\begin{equation}
    \widetilde{E}(\vec{r},\omega-\omega_0)=F(x,y)\widetilde{A}(z,\omega-\omega_0)\exp(i\beta_0z)
\end{equation}
其中，$F(x,y)$ 是光波模式，$\widetilde{A}(z,\omega)$ 是光波的慢变包络函数，$\beta_0$ 是波数。将试探解代入 Helmholtz 方程(\ref{eq:Helmholtz})，得到光波模式方程(\ref{eq:F(x,y)})和慢变包络函数方程(\ref{eq:A(z)})
\begin{subequations}
    \begin{equation}
        \frac{\partial^2F}{\partial x^2}+\frac{\partial^2F}{\partial y^2}+\left[\varepsilon_r(\omega)k_0^2-\widetilde{\beta}^2\right]F=0
        \label{eq:F(x,y)}
    \end{equation}
    \begin{equation}
        2i\beta_0\frac{\partial \widetilde{A}}{\partial z}+\left(\widetilde{\beta}^2-\beta_0^2\right)\widetilde{A}=0
        \label{eq:A(z)}
    \end{equation}
\end{subequations}
包络函数方程(\ref{eq:A(z)})中，需要做慢变包络近似\footnote{行波脉冲的包络在时间和空间上缓慢变化（与波的周期或波长相比）。}，忽略二阶导数 $\partial^2\widetilde{A}/\partial z^2$。

若介质对光场的响应是线性的，那么可用本征函数理论求解(\ref{eq:F(x,y)})。但是，短脉冲在光纤中传输时，介电响应并非线性
\begin{align}
    \varepsilon_r(\omega)&=(n+\Delta n)^2\approx n^2+2n\Delta n\\
    \Delta n(\omega)&=n_2|E|^2+\frac{i\alpha}{2k_0} \nonumber
\end{align}
那么可用微扰理论求解脉冲传输问题\cite{Agrawal}。把 $\Delta n$ 处理为 $n$ 的微扰，先用线性响应 $\varepsilon_r=n^2$ 求解(\ref{eq:F(x,y)})得到模式分部 $F(x,y)$ 和对应的本征值 $\beta$；然后加入微扰 $\Delta n$ 的影响，一阶微扰不影响模式分布 $F(x,y)$，但是本征值修正为
\begin{align}
    \widetilde{\beta}(\omega)&=\beta(\omega)+\Delta\beta(\omega)\\
    \Delta\beta(\omega)&=\frac{\omega^2n(\omega)}{c^2\beta(\omega)}\frac{\iint_{-\infty}^{\infty}\Delta n(\omega)|F(x,y)|^2\mathrm{d}x\mathrm{d}y}{\iint_{-\infty}^{\infty}|F(x,y)|^2\mathrm{d}x\mathrm{d}y} \nonumber
\end{align}

本征值 $\beta(\omega)$ 决定了光波模式，光纤通信领域中它被称为传输常数。传输常数 $\beta$ 主要取决于材料
\begin{equation}
    \beta(\omega)\approx n(\omega)\frac{\omega}{c}
\end{equation}
传输常数 $\beta(\omega)$ 是光波频率的函数，可对其在中心频率 $\omega_0$ 附近做 Taylor 展开以考察其色散效应
\begin{subequations}
    \begin{align}
        \beta(\omega)=\beta_0+&\beta_1(\omega-\omega_0)+\frac{1}{2}\beta_2(\omega-\omega_0)^2+\frac{1}{6}\beta_3(\omega-\omega_0)^3+\cdots \nonumber \\
        &\beta_1=\frac{\mathrm{d}\beta}{\mathrm{d}\omega}=\frac{1}{c}\left(n+\frac{\mathrm{d}n}{\mathrm{d}\omega}\right)=\frac{1}{v_g}\\
        &\beta_2=\frac{1}{c}\left(2\frac{\mathrm{d}n}{\mathrm{d}\omega}+\omega\frac{\mathrm{d}^2n}{\mathrm{d}\omega^2}\right)
    \end{align}
\end{subequations}
其中，参量 $\beta_1$ 与群速度 $v_g$（光脉冲包络移动的速度）有关；参量 $\beta_2$ 描述的就是群速度的色散。群速度色散(group-velocity dispersion，简称 GVD)是造成脉冲展宽的主要原因。

在准单色假设下，中心频率 $\omega_0$ 附近的 $\beta(\omega)$，$\Delta\beta(\omega)$
\begin{subequations}
    \begin{align}
        \beta&=\beta_0+\beta_1(\omega-\omega_0)+\frac{1}{2}\beta_2(\omega-\omega_0)^2\\
        \Delta\beta&=\Delta\beta_0
    \end{align}
\end{subequations}
有了本征值 $\widetilde{\beta}(\omega)$ 的近似值，就可以把包络方程(\ref{eq:A(z)})近似处理为\cite{Agrawal}
\begin{equation}
    \frac{\partial \widetilde{A}}{\partial z}=i[\beta(\omega)+\Delta \beta(\omega)-\beta_0]\widetilde{A}    
\end{equation}
对频域中的包络方程进行 Fourier 逆变换，便得到脉冲包络随时间演化的方程
\begin{equation}
    \frac{\partial A}{\partial z}+\beta_1\frac{\partial A}{\partial t}+\frac{i\beta_2}{2}\frac{\partial^2 A}{\partial t^2}=i\Delta\beta_0A
\end{equation}
根据微扰讨论结果可将脉冲包络在光纤中的传输表达为
\begin{equation}
    \frac{\partial A}{\partial z}+\beta_1\frac{\partial A}{\partial t}+\frac{i\beta_2}{2}\frac{\partial^2 A}{\partial t^2}+\frac{\alpha}{2}A=i\gamma(\omega_0)|A|^2A
\end{equation}
该方程中，$\alpha$ 是光纤的损耗项，$\beta_1$ 描述的是脉冲群速度，$\beta_2$ 是光纤的色散项，$\gamma$ 包含了光纤的非线性效应。

在研究脉冲传输问题时，通常采用以群速度 $v_g$ 随脉冲移动的参考系，这就要求对脉冲传输方程做延时变换\cite{Agrawal}
\begin{equation}
    T=t-z/v_g=t-\beta_1z
\end{equation}
忽略光纤损耗 $\alpha$ 的情况下，脉冲包络传输就用{\bfseries 非线性 Schr\"odinger 方程（noline Schr\"odinger equation，简称 NLSE）}描述\footnote{在研究受激散射和或脉宽小于 1 ps 的超短脉冲传输的问题时，需要更高阶的色散项和非线性项对非线性 Schr\"odinger 方程进行修正，本文不再对推广的非线性 Schr\"odinger 方程进行介绍。}
\begin{equation}
    i\frac{\partial A}{\partial z}=\frac{\beta_2}{2}\frac{\partial^2A}{\partial T^2}-\gamma|A|^2A
    \label{eq:Schrodinger}
\end{equation}

\section{光孤子效应}
\label{sec:optical soliton}
\subsection{孤立波与孤子}
1834 年，英国造船工程师 Russel 骑着马沿着狭窄的运河走时，偶然观察到一个有趣的现象\cite{liangkunmiao}。当小船突然停止前进时，船头会挤出一个水堆。这个水堆保持着自己的形状，以很高的速度沿着河道继续前进。Russel 骑马追了约两英里，它才渐渐消失。Russel 称这种现象为{\bfseries 孤立波（solitary wave）}，并在英国科学协会中报告了这一现象，引起了很长时间的争论。

直到 1895 年，两位荷兰的数学家 Korteweg 和 de Vries 在流体动力学的研究中导出了浅水波方程\cite{liangkunmiao}，即 KdV 方程
\begin{equation}
    \frac{\partial u}{\partial t}+\alpha u\frac{\partial u}{\partial x}+\frac{\partial^3 u}{\partial x^3}=0
\end{equation}
同时，他们还得到这一方程的一种行波解，在波长趋于无限长的情形下，该行波解能很好地描述 Russel 所发现的孤立波。至此，人们才渐渐接受了 Russel 提出的孤立波的概念，但是这项研究并没有得到人们的广泛关注。

1955 年，Fermi、Pasta 和 Ulam 在一篇有关晶格热传导的论文中提出了一个看似与孤立波无关的问题 —— 在具有非线性相互作用的一维振子链模型中，如果按照能量均分定理，初始能量最终会均分到振子链的所有自由度。但是所得结果却是经过一段时间，能量又基本回到最初的振动模式（FPU 问题）\cite{liangkunmiao}。

一直到 1965 年，美国普林斯顿大学的 Kruskal 和贝尔实验室的 Zabusky 对 FPU 问题进行数值计算。他们发现，在一环形振子链\footnote{周期性边界条件}中，初始脉冲随着时间的推移，演化成几个大小不同的孤立波，它们互相碰撞、穿行，碰撞后仍然保持各自原来的形状与速度，并在某一时刻，所有孤立波在环上同一地点相遇，基本重现初始状态。由于这种类似粒子的性质，Kruskal 和 Zabusky 将这样的孤立波命名为{\bfseries 孤子（soliton）}。{\bfseries 孤子是一种特殊的孤立波，总结其特点\cite{yangbojun,liangkunmiao}
\begin{enumerate}[label=(\arabic*)]
    \item 孤子之间碰撞和相互作用后，能保持其形状和速度不变，仅相位发生变化；
    \item 孤子的能量总是集中在有限的空间区域内，不会耗散到无限的空间区域中去。
\end{enumerate}}

\subsection{孤子本质}
\label{sec:soliton}
目前为止，人们在很多学科领域发现孤子现象，并且在理论上得到了一批描述孤子的非线性偏微分方程（见图\ref{fig:soliton-eq}），孤子问题也逐渐发展为完整、系统的理论。虽然孤子是非线性偏微分方程的一种特殊解，但是孤子的物理本质并不复杂，它是色散效应和非线性效应综合作用的结果。
\begin{figure}[tbp]
    \centering
    \includegraphics[width=0.9\textwidth]{孤子方程.png}
    \caption{不同领域中的孤子问题}
    \label{fig:soliton-eq}
\end{figure}

不同频率波的传播速度不同形成色散，色散使脉冲（波包）在传输过程中逐渐展宽，能量逐渐耗散，最终导致脉冲（波包）消失；非线性波的传播速度依赖于振幅，振幅越大的地方传播速度越快，这就使得波的前沿不断变陡，波形更加尖锐。{\bfseries 因此，色散作用使波形平展，非线性效应使波形尖锐，两者的作用刚好相反。当两种因素相互平衡时，便能使波形保持不变，形成孤子\cite{liangkunmiao}。}

如\ref{sec:Schrodinger}节所述，非线性 Schr\"odinger 方程描述的就是色散效应和非线性效应共同作用下脉冲在光纤中的传输过程。Hasegawa 和 Tappert 就在 1973 年指出，反常群速度色散和非线性效应的相互作用使得孤子脉冲可以在低损光纤中传输，并通过数值实验证实了光孤子在通信系统中的稳定性\cite{Hasegawa}。

为了描述到底是色散效应还是非线性效应对脉冲的传输演化起主要作用，需要对非线性 Schr\"odinger 方程进行无量纲化处理，并引入色散长度 $L_D$ 和非线性长度 $L_{NL}$ 来衡量脉冲传输过程中色散和非线性效应哪个起重要作用。那么，先对脉冲进行归一化处理
\begin{align}
    &\tau=\frac{T}{T_0}=\frac{t-z/v_g}{T_0} \nonumber \\
    &A(z,\tau)=\sqrt{P_0}U(z,\tau)
\end{align}
$T_0$ 为初始脉宽，$\tau$ 是归一化时间尺度（不具有时间量纲），$P_0$ 为入射脉冲峰值功率，$U$ 是归一化振幅（也不具备振幅量纲），已忽略光纤的损耗系数 $\alpha$。根据非线性 Schr\"odinger 方程(\ref{eq:Schrodinger})， 归一化振幅 $U(z,\tau)$ 需满足
\begin{align}
    i\frac{\partial U}{\partial z}&=\frac{\mathrm{sng}(\beta_2)}{2L_D}\frac{\partial^2U}{\partial \tau^2}-\frac{1}{L_{NL}}|U|^2U\\
    L_D&=\frac{T^2_0}{|\beta_2|}\hspace{4ex}L_{NL}=\frac{1}{\gamma P_0} \nonumber
\end{align}

可根据脉冲传输距离 $z$ 与色散长度 $L_D$ 和非线性长度 $L_{NL}$ 之间的相对大小，将光纤分为不同的传输区域\cite{Agrawal}
\begin{enumerate}[label=(\arabic*)]
    \item $z\ll L_{D}$ 且 $z\ll L_{NL}$，色散和非线性效应都不起重要作用。传输过程中，脉冲形状保持不变（见图\ref{fig:distortionless}）。
    \item $z\ll L_{NL}$ 而 $z\approx L_{D}$，脉冲传输演化过程中，群速度色散（GVD）起主要作用，使脉冲展宽（见图\ref{fig:GVD}）。
    \item $z\ll L_{D}$ 但 $z\approx L_{NL}$，非线性效应对脉冲起主要作用。其结果是产生自相位调制（self-phase modulation，SPM）现象，使脉冲频谱发生变化。较强的自相位调制和较弱的群速度色散导致光波分裂现象，脉冲前后沿变陡，并演化出快速振荡的精细结构（见图\ref{fig:SPM_time}），频谱中央也因此产生多峰结构\cite{Agrawal}（见图\ref{fig:SPM_frequent}）。
    \item $z\approx L_{D}$，$z\approx L_{NL}$，色散和非线性效应对脉冲传输都有重要作用。在正常色散区 $\beta_2>0$，群速度色散和自相位调制可用于脉冲压缩；在反常色散区 $\beta_2<0$，两者的作用能形成光孤子\cite{Agrawal}。
\end{enumerate}
\begin{figure}[tbp]
    \centering
    \begin{minipage}[t]{0.45\linewidth}
        \centering
        \includegraphics[width=1\linewidth]{脉冲无畸变传输.png}
        \caption{脉冲无畸变传输}
        \label{fig:distortionless}
    \end{minipage}       
    \begin{minipage}[t]{0.45\linewidth}
        \centering
        \includegraphics[width=1\linewidth]{脉冲展宽.png}
        \caption{群速度色散导致脉冲展宽}
        \label{fig:GVD}    
    \end{minipage}
\end{figure}
\begin{figure}[tbp]
    \centering
    \begin{subfigure}[t]{0.45\linewidth}
        \begin{minipage}[b]{1\linewidth}
        \includegraphics[width=1\linewidth]{自相位调制_时域.png}
        \caption{时域}
        \label{fig:SPM_time}
        \end{minipage}
    \end{subfigure}
    \begin{subfigure}[t]{0.45\linewidth}
        \begin{minipage}[b]{1\linewidth}
        \includegraphics[width=1\linewidth]{自相位调制_频域.png}
        \caption{频域}
        \label{fig:SPM_frequent}
        \end{minipage}
    \end{subfigure}
    \caption{自相位调制主要影响下的光波分裂}
\end{figure}

\subsection{逆散射方法}
如 \ref{sec:soliton} 节所述，孤子现象由一些特殊的非线性偏微分方程所描述，人们在研究这些特殊的非线性偏微分方程的过程中发展了量子力学中的逆散射方法（the Inverse Scattering Transform Method，简称 IST）。

1967 年，Gardner 等人将 IST 方法应用于 KdV 方程求解，对孤子现象做出严格的数学证明。一开始，人们认为 IST 方法能用于求解 KdV 方程仅仅是因为 Schr\"odinger 方程与 KdV 方程之间的巧妙联系。但是，Lax 在 1968 年对 IST 方法做了推广，使得 IST 方法能求解更多的非线性偏微分方程。1971 年，Zakharov 和 Shabat 将 IST 方法用于非线性 Schr\"odinger 方程的求解，并给出了非线性 Schr\"odinger 方程的解析解\cite{yangbojun,Zakharov}。

量子力学中，Schr\"dinger 方程的散射问题是根据已知势场求解散射数据以及波函数在无穷远处的渐近行为；逆散射问题正好相反，已知散射数据和波函数在无穷远处的边界条件反向求解散射势\footnote{逆散射方法最初就是利用散射数据来构造势能曲线，以了解原子与分子结构。}。{\bfseries IST 方法并不是直接去求非线性偏微分方程的解，而是把非线性演化问题作为一个散射问题考虑，根据初始散射数据求解散射数据在散射过程中的演化，再由演化的散射数据重构位势，散射势就是非线性偏微分方程的解。IST 方法的思想与 Fourier 变换法求解线性偏微分方程非常类似（见图\ref{fig:Fourier}），因此 IST 方法也常被称为非线性 Fourier 变换法\cite{liangkunmiao}。}
\begin{figure}[tbp]
    \centering
    \begin{subfigure}[t]{0.40\linewidth}
        \begin{minipage}[b]{1\linewidth}
        \includegraphics[width=1\linewidth]{线性傅里叶变换.png}
        \caption{线性Fourier变换法}
        \label{fig:Fourier_pde}
        \end{minipage}
    \end{subfigure}
    \hspace{4ex}
    \begin{subfigure}[t]{0.40\linewidth}
        \begin{minipage}[b]{1\linewidth}
        \includegraphics[width=1\linewidth]{非线性傅里叶变换.png}
        \caption{非线性Fourier变换法}
        \end{minipage}
    \end{subfigure}
    \caption{Fourier变换方法与逆散射方法}
    \label{fig:Fourier}
\end{figure}

本节着重讨论非线性 Schr\"odinger 方程的 IST 方法求解。不妨对传输距离 $z$ 也做无量纲化处理 $\xi=z/L_D$，并引入参量 $N^2=L_{D}/L_{NL}$，把非线性 Schr\"odinger 表述成更简洁的形式
\begin{equation}
    i\frac{\partial U}{\partial \xi}=\mathrm{sng}(\beta_2)\frac{1}{2}\frac{\partial^2 U}{\partial \tau^2}-N^2|U|^2U
\end{equation}
在反常色散区 $\beta_2<0$，标准化的非线性 Schr\"odinger 方程为
\begin{align}
    &u=NU=\sqrt{\gamma L_D} A\nonumber \\
    &i\frac{\partial u}{\partial \xi}+\frac{1}{2}\frac{\partial^2 u}{\partial \tau^2}+|u|^2u=0
    \label{eq:NLSE}
\end{align}
与该非线性演化相联系的散射问题是\cite{guoyucui}
\begin{equation}
    \begin{pmatrix}
        i\partial/\partial \tau & u \\
        -u^* & -i\partial/\partial \tau
    \end{pmatrix}\begin{bmatrix}
        \psi_1\\
        \psi_2
    \end{bmatrix}=k\begin{bmatrix}
        \psi_1\\
        \psi_2
    \end{bmatrix}
    \label{eq:scatter}
\end{equation}
其中，$\psi_1$ 和 $\psi_2$ 是被势场 $u(\xi,\tau)$ 散射的两个波函数，$k$ 是本征值。根据 Lax 推广的 IST 方法，非线性偏微分方程应等价为两个线性问题，其中一个就是如(\ref{eq:scatter})所示的散射问题；另一个线性方程描述波函数随空间传输的演化过程
\begin{equation}
    \frac{\partial}{\partial \xi}\begin{bmatrix}
        \psi_1\\
        \psi_2
    \end{bmatrix}=\textbf{\textit{M}}\begin{bmatrix}
        \psi_1\\
        \psi_2
    \end{bmatrix}
    \label{eq:wave function evolution}
\end{equation}
算子 $\textbf{\textit{M}}$ 的具体形式需要用 AKNS 方法\footnote{1973 年，Ablowitz、Kaup、Newell 和 Segur 在求解 Sine-Gordon 方程时提出由散射问题算子 $\hat{L}$ 构造波函数演化算子 $\hat{M}$ 的方法}推导\cite{guoyucui}，结果比较繁琐，此处不再说明。

散射问题(\ref{eq:scatter})是一个本征值问题。{\bfseries 当本征值 $k$ 为实数时，可以取连续值，本征函数为行波解，即波函数对应于连续态}，其在无穷远处的渐进行为
\begin{equation}
    \begin{bmatrix}
        \psi_1\\
        \psi_2
    \end{bmatrix}=\left\{\begin{aligned}
        &\begin{bmatrix}
            1\\
            0
        \end{bmatrix}e^{-ik\tau}+R(k,\xi)\begin{bmatrix}
            0\\
            1
        \end{bmatrix}e^{ik\tau}, & \tau\to-\infty\\
        &T(k,\xi)\begin{bmatrix}
            1\\
            0
        \end{bmatrix}e^{-ik\xi}, & \tau\to\infty
    \end{aligned}\right.
\end{equation}
其中，$R(k,\xi)$ 为反射系数，$T(k,\xi)$ 为透射系数。{\bfseries 当本征值 $k$ 为复数时，只能取有限个分立值 $k_n$，本征函数指数衰减，即波函数对应为束缚态（孤子）}。这种情形下，本征值 $k_n$ 正是反射系数 $R(k,\xi)$ 在 $k$ 平面上的极点，对应的留数为 $C_{n}(k_n,\xi)$\cite{guoyucui}。

连续本征值 $k$ 对应的反射系数 $R(k,\xi)$ 和离散本征值 $k_n$ 对应的留数 $C(k_n,\xi)$ 就是散射问题(\ref{eq:scatter})的散射数据。可由波函数演化方程(\ref{eq:wave function evolution})确定散射数据随传输距离的演化关系
\begin{subequations}
    \begin{align}
        R(k,\xi)&=R(k,0)e^{i2k^2\xi}\\
        C_n(k_n,\xi)&=C_n(k_n,0)e^{ik_n^2\xi}
    \end{align}
\end{subequations}


根据 $\xi=0$ 处的输入脉冲得到初始的散射数据 $R(k,0)$ 和 $C_n(k_n,0)$\footnote{1974 年，Satsuma 和 Yajima 讨论了初始条件对非线性 Schr\"odinger 方程解的影响。}，脉冲在光纤中的传输对应于散射数据演化，直接给出逆散射问题的结果，也就是散射势
\begin{equation}
    u(\xi,\tau)=-2K_1(\tau,\tau,\xi)
\end{equation}
其中，函数 $K_1(x,y,t)$ 满足 Gel'fand-Levitan-Marchenko 积分方程（简称 GLM 方程）
\begin{equation}
    \left\{\begin{aligned}
        &K_1(x,y,t)=B^*(x+y,t)+\int_x^\infty B^*(y+z,t)K_2^*(x,z,t)\mathrm{d}z\\
        &K_2^*(x,y,t)=-\int_x^\infty B(y+z,t)K_1(x,z)\mathrm{d}z
    \end{aligned}\right.
    \label{eq:GLM}
\end{equation}
积分核 $B(x+y,\xi)$ 由散射数据给出
\begin{equation}
    B(x+y,\xi)=\sum_{n=1}^NC^2_n(ik_n,\xi)e^{k_n(x+y)}+\frac{1}{2\pi}\int_{-\infty}^{\infty}R(k,\xi)e^{ik(x+y)}\mathrm{d}k
\end{equation}

上述过程给出了严格求解非线性 Schrodinger 方程的步骤，其中不可避免地要求解复杂的积分方程(\ref{eq:GLM})。但是，对于无反射的孤子解情形 $R(k,0)=0$，便可通过求解代数方程组得到 $u(\xi,\tau)$
\begin{align}
    u(\xi,\tau)&=-2\sum_{n=1}^{N}\lambda_n^*\psi_{2n}^* 
    \label{eq:soliton}\\
    \lambda_n^*&=\sqrt{C_n}e^{i(k_n\tau+k_n^2\xi)} \nonumber
\end{align}
其中，$\psi_{2n}^*$ 满足代数方程组
\begin{equation}
    \begin{aligned}
        \psi_{1n}+\sum_{n'=1}^N\frac{\lambda_n\lambda_{n'}^*}{k_n-k^*_{n'}}\psi^*_{2n'}&=0\\
        \psi_{2n}^*-\sum_{n'=1}^N\frac{\lambda^*_{n}\lambda_{n'}}{k^*_n-k_{n'}}\psi_{1n'}&=\lambda_n^*
    \end{aligned}
\end{equation}
极点数目 $N$ 或离散本征值 $k_n$ 的数目表征孤子阶数。
\subsection{光孤子}
从 IST 方法求解非线性 Schr\"odinger 方程的过程我们清楚地可以知道，光孤子解来自散射问题的束缚态（$R(k,\xi)$ 的极点处，本征值为复数 $k_n$）。上一节仅在反常群速度色散区讨论光孤子的解，随着光纤和激光技术的成熟，人们在实验上发现或者在理论上提出了多种光孤子（见图\ref{fig:optical soliton}），本文重点讨论反常色散区的光孤子（亮孤子）\cite{Agrawal}。
\begin{figure}[hbtp]
    \centering
    \includegraphics[width=0.6\textwidth]{光孤子.png}
    \caption{不同类型的光孤子}
    \label{fig:optical soliton}
\end{figure}
\subsubsection{基阶孤子}
单个本征值 $N=1$ 的情形下，脉冲传输能形成基阶孤子，可通过求解孤子解方程组(\ref{eq:soliton})得到基阶孤子的一般形式
\begin{equation}
    u(\xi,\tau)=\eta\mathrm{sech}[\eta(\tau-\tau_s+\delta\xi)]\exp[i(\eta^2-\delta^2)\xi/2-i\delta\tau+i\phi_s]
\end{equation}
其中，$\eta$、$\delta$、$\tau_s$ 和 $\phi_s$ 是表征基阶孤子的 4 个参数。$\phi_s$ 是相位常数，在研究单孤子问题时常把它忽略，在讨论孤子相互作用问题时才起作用。参量 $\tau_s$ 表示孤子峰的位置，选取适当的时间原点可以使 $\tau_s=0$。参量 $\eta$ 不仅决定孤子的振幅，而且同时决定了孤子的宽度。参量 $\delta$ 表示孤子相对于中心频率 $\omega_0$ 的频移，同时也导致群速度 $v_g$ 的改变\cite{Agrawal}
\begin{equation}
    \begin{aligned}
        |u(\xi,\tau)|&=\eta\mathrm{sech}[\eta(\tau+\delta\xi)]=\eta\mathrm{sech}\left[\eta\left(\frac{t-\beta_1z}{T_0}+\delta\frac{z}{L_D}\right)\right]\\
        &=\eta\mathrm{sech}\left(\eta\frac{t-\beta'_1z}{T_0}\right)\\
        \beta'_1&=\beta_1-\frac{\delta|\beta_2|}{T_0}
    \end{aligned}
\end{equation}

在忽略相位常量 $\phi_s$、孤子中心 $\tau_s$ 和频移 $\delta$ 参量后，基阶孤子可以描述为
\begin{equation}
    u(\xi,\tau)=\eta\mathrm{sech}(\eta\tau)\exp\left(i\frac{\eta^2\xi}{2}\right)
\end{equation}
从中我们可以发现
\begin{enumerate}[label=(\arabic*)]
    \item 孤子宽度 $T_0/\eta$ 反比于孤子振幅 $\eta$，这是孤子很重要的一个特点；
    \item 对于 $N=1$ 的基阶孤子脉冲，在光纤中将无畸变地传输（见图\ref{fig:first soliton}），这使得基阶孤子在光纤通信中具有广阔前景。
\end{enumerate}
\begin{figure}[bp]
    \centering
    \begin{minipage}[t]{0.45\linewidth}
        \centering
        \includegraphics[width=1\linewidth]{基阶孤子无畸变传输.png}
        \caption{基阶孤子无畸变传输}
        \label{fig:first soliton}
    \end{minipage}       
    \begin{minipage}[t]{0.45\linewidth}
        \centering
        \includegraphics[width=1\linewidth]{二阶孤子解析解.png}
        \caption{二阶孤子的周期性演化}
        \label{fig:second soliton}    
    \end{minipage}
\end{figure}
\subsubsection{二阶和高阶孤子}
高阶孤子也可以用通解(\ref{eq:soliton})来描述，不同的本征值 $k_n$ 和留数 $C_n$ 可以构造出各种各样的高阶孤子。可根据本征值与留数的关系区分不同的孤子集，其中一类孤子集其初始形状如下
\begin{equation}
    u(0,\tau)=N\mathrm{sech}(\tau)
\end{equation}
孤子阶数 $N$ 为整数\footnote{当 $N$ 不为整数时，涉及到孤子稳定性问题，本文不做讨论。}。这一类孤子有一个很有趣的特征——包络形状 $|u(\xi,\tau)|$ 做周期 $\xi_0=\pi/2$ 的周期性变化。拿二阶孤子举例，对于本征值 $k_1=1/2$ 和 $k_2=3i/2$ 的二阶孤子，其场分布 $u(\xi,\tau)$
\begin{equation}
    u(\xi,\tau)=\frac{4[\cosh(3\tau)+3e^{4i\xi}\cosh(\tau)]}{\cosh(4\tau)+4\cosh(2\tau)+3\cos(4\xi)}e^{i\xi/2}
\end{equation}
根据场分布 $u(\xi,\tau)$ 的解析表达式作图（见图\ref{fig:second soliton}），可以发现孤子周期 $z_0$
\begin{equation}
    z_0=\xi_0L_D=\frac{\pi}{2}\frac{T^2_0}{|\beta_2|^2}
\end{equation}

对三阶孤子（$N=3$）进行数值模拟（见图\ref{fig:third soliton_time}）。发现脉冲在光纤中传输时，开始阶段脉宽被压缩，在 $z_0/2$ 处被窄化的脉冲一分为二，在孤子周期 $z_0$ 处又变回原来形状的单个脉冲。脉冲变窄、分裂、再恢复的过程在每个孤子周期内重复进行。

从三阶孤子的频谱演化过程来看（见图\ref{fig:third soliton_frequent}），在 $z=0.3L_D$ 附近，出现了自相位调制的典型振荡结构，这种振荡性的频谱结构使得孤子的前沿和后沿发生变化，导致脉冲分裂；随后的反常群速度色散导致脉宽变窄，脉冲频谱也相应改变。也就是说，自相位调制和群速度色散交替的主导作用使得光孤子无论是时域还是频域上来看都将呈现周期性的演化过程\cite{Agrawal}。
\begin{figure}[btp]
    \centering
    \begin{subfigure}[t]{0.45\linewidth}
        \begin{minipage}[b]{1\linewidth}
        \includegraphics[width=1\linewidth]{三阶孤子演化_时域.png}
        \caption{时域}
        \label{fig:third soliton_time}
        \end{minipage}
    \end{subfigure}
    \begin{subfigure}[t]{0.45\linewidth}
        \begin{minipage}[b]{1\linewidth}
        \includegraphics[width=1\linewidth]{三阶孤子演化_频域.png}
        \caption{频域}
        \label{fig:third soliton_frequent}    
        \end{minipage}
    \end{subfigure}
    \caption{三阶孤子的周期性演化}
\end{figure}

{\bfseries 对于基阶孤子（$N=1$），自相位调制和色散影响互相平衡，脉冲的形状和频谱几乎不随传输变化。对于二阶、三阶甚至更高阶的孤子，自相位调制和群速度色散的影响地位是交替变化着的，导致传输过程中出现脉冲变窄、分裂、再恢复的周期性演化结果\cite{Agrawal}。}

\section{数值方法}
\label{sec:numerical}
数学上可以用逆散射方法得到非线性 Schr\"odinger 方程的解析解，也可以用矩方法、变分法得到半解析半近似的解。虽然解析解能很好地反映脉冲在光纤中传输的物理图像，但是大部分实际问题求解析解非常繁琐或是根本就无法用解析方法求解。因此通过数值方法求解非线性 Schr\"odinger 方程来研究光脉冲的非线性传输是必要的。求解这种非线性偏微分方程的数值方法大致可以分为两类
\begin{enumerate}[label=(\arabic*)]
    \item 通过频谱间接求解，分步 Fourier 方法（split-step Fourier method，简称 SSFM）和伪谱方法（pseudospectral method）；
    \item 在时域中直接求解，有限差分法（finite difference method，简称 FDM）。
\end{enumerate}
\subsection{伪谱方法}
{\bfseries Fourier 积分变换可以把时域（或空间域）的偏微分方程转换为频域的常微分方程，通过求解频域的常微分方程，再进行逆 Fourier 变换得到偏微分方程的时域（空间域）的解（见图\ref{fig:Fourier_pde}）。将此方法应用于偏微分方程离散求解，就是伪谱方法\cite{zhangxiao}。}

要在计算机上实现伪谱方法，有两个问题需要解决
\begin{enumerate}[label=(\arabic*)]
    \item 实现离散数据序列时域与频域的互相转换；
    \item 实现常微分方程的数值求解。
\end{enumerate}
快速 Fourier 变换（fast Fourier transform，简称 FFT）可以很方便地实现离散 Fourier 变换（MATLAB 提供了快速 Fourier 变换的指令 fft 和 ifft）；通常用 Runge-Kutta 方法求解常微分方程的数值解（MATLAB 中集成了非常多用 Runge-Kutta 方法设计的 ODE 求解器，比如最常用的 ode45）\cite{Anne}。因此，应用伪谱方法求解偏微分方程在技术上是可行的，求解步骤也十分清楚\cite{zhangxiao}
\begin{enumerate}[label=(\arabic*)]
    \item 对偏微分方程做时域（或空间域）的 Fourier 变换，再将其离散化；
    \item 调用 ODE 求解器对离散化的常微分方程进行求解；
    \item 用 ifft 函数将频域上的计算结果转换为时域（或空间域）上的结果。
\end{enumerate}

1978 年，Fornberg 和 Whitham 提出用伪谱方法求解非线性
Schrödinger 方程\cite{Fornberg}。对于非线性 Schr\"odinger 方程(\ref{eq:NLSE})，通常选择做时域 Fourier 变换，得到脉冲传输（频域）的常微分方程\footnote{应用伪谱方法求解非线性 Schr\"odinger 方程要求脉冲在时域具有周期性，这在光纤传输中是可行的。}
\begin{align}
    &\frac{\partial \widetilde{u}}{\partial \xi}=-\frac{i\omega^2}{2}\widetilde{u}+i\mathscr{F}\left[\left|\mathscr{F}^{-1}[\widetilde{u}]\right|^2\mathscr{F}^{-1}[\widetilde{u}]\right]\\
    &\widetilde{u}(0,\omega)=\mathscr{F}\left[N\cdot\mathrm{sech}(\tau)\right]
\end{align}
将该方程离散化后调用 ode45 指令求解离散后的频域演化方程，再用 ifft 函数将频域结果转换为时域演化过程（详见附录1）。

图\ref{fig:third soliton_time} 和图\ref{fig:third soliton_frequent} 就是用伪谱方法计算输入脉冲 $u(0,\tau)=3\mathrm{sech(\tau)}$ 在光纤中传输得到的三阶光孤子演化结果。用伪谱方法对高阶孤子进行数值模拟，同样得到了脉冲窄化、分裂、再恢复的周期性物理过程（见图\ref{fig:Fourier spectral method}）。
\begin{figure}[tbp]
    \centering
    \begin{subfigure}[t]{0.85\linewidth}
        \captionsetup{justification=centering} 
        \begin{minipage}[b]{1\linewidth}
        \includegraphics[width=0.45\linewidth]{FSM_四阶孤子演化_侧视图.png}
        \includegraphics[width=0.45\linewidth]{FSM_四阶孤子演化_俯视图.png}
        \caption{$u(0,\tau)=4\mathrm{sech}(\tau)$}
        \end{minipage}
    \end{subfigure}\\
    \begin{subfigure}[t]{0.85\linewidth}
        \captionsetup{justification=centering} 
        \begin{minipage}[b]{1\linewidth}
        \includegraphics[width=0.45\linewidth]{FSM_五阶孤子演化_侧视图.png}
        \includegraphics[width=0.45\linewidth]{FSM_五阶孤子演化_俯视图.png}
        \caption{$u(0,\tau)=5\mathrm{sech}(\tau)$}
        \end{minipage}
    \end{subfigure}
    \caption{伪谱方法模拟的高阶孤子周期性演化}
    \label{fig:Fourier spectral method}
\end{figure}
\subsection{分步 Fourier 方法}
伪谱方法在计算机上求解效率高、计算内存小，但是它需要对频域的常微分方程进行数值求解。并且在研究受激 Raman 散射和 受激 Brillouin 散射问题时，需要在非线性 Schr\"odinger 方程中引入高阶色散项和非线性项，这种情形下需要对频域的常微分方程进行降阶等一系列处理，伪谱方法的可拓展性就非常有限。其实早在伪谱方法之前，人们在求解 KdV 方程的时候就应用了分步 Fourier 方法\cite{Tappert,Chellappan}，后来这一数值方法被推广到非线性 Schr\"odinger 方程及其他非线性偏微分方程的求解。
\subsubsection{交替计算}
脉冲在光纤中传输时，光场包络同时受色散和非线性效应的影响
\begin{align}
    \frac{\partial A}{\partial z}&=\left(\hat{D}+\hat{N}\right)A \\
    \hat{D}&=-\frac{\alpha}{2}-\frac{i\beta_2}{2}\frac{\partial^2}{\partial T^2}+\cdots \nonumber \\
    \hat{N}&=i\gamma|A|^2+\cdots \nonumber
\end{align}
也就是说，色散效应（线性微分算子 $\hat{D}$）和非线性效应（非线性函数 $\hat{N}$）对包络函数 $A$ 的作用是同步的。分步 Fourier 方法假设光场传输过程中这两种作用交替进行，分别考虑其影响。这种假设从计算结果上看是可行的，能得到光波传输的近似结果\cite{Agrawal,Fisher}。

脉冲从 $z$ 到 $z+h$ 传播一小段距离，这一小段传输过程可以分两步进行：第一步，仅考虑非线性效应的影响，忽略线性算子的作用，$\hat{D}=0$；第二步，仅考虑色散效应的作用，忽略非线性函数的影响，$\hat{N}=0$
\begin{equation}
    A(z+h,T)\approx\exp(h\hat{D})\exp(h\hat{N})A(z,T)
\end{equation}
非线性效应的影响 $\exp(h\hat{N})$ 可以在时域中计算，色散效应的影响 $\exp(h\hat{D})$ 在频域中计算更为方便
\begin{equation}
    \exp(h\hat{D})B(z,T)=\mathscr{F}^{-1}\left[\exp(h\hat{D}(-i\omega))\mathscr{F}[B(z,T)]\right]
\end{equation}
其中，$\hat{D}(-i\omega)$ 表示用 $-i\omega$ 去替换时域微分算子 $\partial/\partial T$\cite{Agrawal}。

最原始的分步 Fourier 方法是非常不精确的，因为它忽略了算子 $\hat{D}$ 和算子 $\hat{N}$ 的非对易性，根据 Baker-Campbell-Hausdorf（BCH）公式
\begin{equation}
    \exp(\hat{a})\exp(\hat{b})=\exp\left(\hat{a}+\hat{b}+\frac{1}{2}[\hat{a},\hat{b}]+\frac{1}{12}[\hat{a}-\hat{b},[\hat{a},\hat{b}]]+\cdots\right)
\end{equation}
可知对易子 $h^2[\hat{D},\hat{N}]$ 是分步 Fourier 方法误差的主要来源，因此分步 Fourier 方法产生的空间截断误差是 $O(h^2)$。
\subsubsection{改进分步 Fourier 方法}
可以对脉冲在 $z$ 到 $z+h$ 小区间上传输过程做不同的假设，以提高分步 Fourier 方法的精度。比较常用的一种方案是对称分步 Fourier 方法，该方法假定非线性效应作用在每个小区间的中间
\begin{equation}
    A(z+h,T)\approx\exp\left(\frac{h}{2}\hat{D}\right)\exp\left(\int_z^{z+h}\hat{N}(z')\mathrm{d}z'\right)\exp\left(\frac{h}{2}\hat{D}\right)A(z,T)
    \label{eq:SSFM}
\end{equation}
{\bfseries 分步 Fourier 方法的物理过程是非常直观的（见图\ref{fig:SSFM}）。脉冲 $A(z,T)$ 在最初的 $h/2$ 距离上传输时仅考虑色散作用；在 $z+h/2$ 的处，包络乘以非线性的积分项，表示整个小区间上的非线性效应；脉冲在剩下的 $h/2$ 距离上的传输也只与色散有关，最终得到 $A(z+h,T)$\cite{Agrawal}。}
\begin{figure}[btp]
    \centering
    \includegraphics[width=0.5\textwidth]{SSFM.png}
    \caption{分步 Fourier 方法示意图}
    \label{fig:SSFM}
\end{figure}

应用 BCH 公式可知对称分步 Fourier 方法的主要误差来自双对易子 $h^3[\hat{N},[\hat{D},\hat{N}]]$ 和 $h^3[\hat{D},[\hat{D},\hat{N}]]$，因此分步 Fourier 方法的截断误差 $O(h^3)$。

分步 Fourier 方法很容易在计算机上实现。可将长度为 $L$ 的光纤划分为 $M$ 个计算区间，若步长 $h=L/M$ 足够小，包含非线性算子积分的指数项可以近似为 $\exp(h\hat{N})$。应用(\ref{eq:SSFM})给出的传输假设
\begin{equation}
    A(L,T)\approx e^{-\frac{1}{2}h\hat{D}}\left(\prod_{m=1}^Me^{h\hat{D}}e^{h\hat{N}}\right)e^{\frac{1}{2}h\hat{D}}A(0,T)
\end{equation}
当然也可以把色散作用集中在每一个小区间的中间，得到另一种不同的算法
\begin{equation}
    A(L,T)\approx e^{-\frac{1}{2}h\hat{N}}\left(\prod_{m=1}^Me^{h\hat{N}}e^{h\hat{D}}\right)e^{\frac{1}{2}h\hat{N}}A(0,T)
    \label{eq:1/2N->D->1/2N}
\end{equation}

图\ref{fig:SSFM_second_soliton} 给出了用对称分步 Fourier 方法(\ref{eq:1/2N->D->1/2N})实现的二阶孤子的时域和频域演化过程（计算代码详见附录2）；图\ref{fig:FSM_second_soliton} 给出的是相同步长 $h$ 下伪谱方法计算所得的结果。将其与二阶孤子解析结果（见图\ref{fig:second soliton}）进行对比发现
\begin{enumerate}[label=(\arabic*)]
    \item 伪谱方法脉冲传输的计算结果与解析解非常吻合；而分步 Fourier 方法却不能很好地体现孤子周期性，并且随着传输距离的增加，计算结果会在在脉冲的前沿和后沿出现振荡；
    \item 从频域结果看，伪谱方法所得频谱可以体现自相位调制和反常色散交替影响的结果；而分步 Fourier 方法所得结果中出现了频移的小峰，这部分较小的频移导致脉冲边缘在传输过程中的振荡。
\end{enumerate}
\begin{figure}[tbp]
    \centering
    \begin{subfigure}[t]{0.85\linewidth}
        \captionsetup{justification=centering} 
        \begin{minipage}[b]{1\linewidth}
        \includegraphics[width=0.45\linewidth]{SSFM_二阶孤子_时域.png}
        \includegraphics[width=0.45\linewidth]{SSFM_二阶孤子_频域.png}
        \caption{分步 Fourier 方法求解孤子时域演化和频谱（$h=0.2$）}
        \label{fig:SSFM_second_soliton}
        \end{minipage}
    \end{subfigure}\\
    \begin{subfigure}[t]{0.85\linewidth}
        \captionsetup{justification=centering} 
        \begin{minipage}[b]{1\linewidth}
        \includegraphics[width=0.45\linewidth]{FSM_二阶孤子_时域.png}
        \includegraphics[width=0.45\linewidth]{FSM_二阶孤子_频域.png}
        \caption{伪谱方法求解孤子时域演化和频谱（$h=0.2$）}
        \label{fig:FSM_second_soliton}
        \end{minipage}
    \end{subfigure}
    \caption{分步 Fourier 方法和伪谱方法求解二阶孤子结果对比}
\end{figure}
导致分步 Fourier 方法计算结果不准确的主要原因是用 $\exp(h\hat{N})$ 代替了非线性算子积分的指数项进行求解。可用迭代的方法计算近似积分，过程比较繁琐，但是通过减小区间步长 $h$ 一样可以得到更可靠的结果。因此，采用更小的空间步长计算，分步 Fourier 方法可以得到和伪谱方法类似的求解结果（见图\ref{fig:split-step Fourier method}）。
\begin{figure}[tbp]
    \centering
    \begin{subfigure}[t]{0.85\linewidth}
        \captionsetup{justification=centering} 
        \begin{minipage}[b]{1\linewidth}
        \includegraphics[width=0.45\linewidth]{SSFM_四阶孤子_侧视图.png}
        \includegraphics[width=0.45\linewidth]{SSFM_四阶孤子_俯视图.png}
        \caption{$u(0,\tau)=4\mathrm{sech}(\tau)$}
        \end{minipage}
    \end{subfigure}\\
    \begin{subfigure}[t]{0.85\linewidth}
        \captionsetup{justification=centering} 
        \begin{minipage}[b]{1\linewidth}
        \includegraphics[width=0.45\linewidth]{SSFM_五阶孤子_侧视图.png}
        \includegraphics[width=0.45\linewidth]{SSFM_五阶孤子_俯视图.png}
        \caption{$u(0,\tau)=5\mathrm{sech}(\tau)$}
        \end{minipage}
    \end{subfigure}
    \caption{分步 Fourier 方法模拟的高阶孤子周期性演化}
    \label{fig:split-step Fourier method}
\end{figure}

\subsection{有限差分法}
{\bfseries 有限差分法是人们在研究微分方程时发展起来的一种数值技术，该方法的基本思想是用差分近似代替微分方程中的微分，把复杂的微分方程问题转化为计算机上容易处理的差分方程（代数方程）问题。利用有限差分法求解偏微分方程基本包括三个步骤\cite{yangbojun}
\begin{enumerate}[label=(\arabic*)]
    \item 将求解区域离散化（划分网格节点）；
    \item 对偏微分方程和定解条件做差分近似（构造差分格式）；
    \item 求解差分方程组（一般采用迭代法求解）。
\end{enumerate}}

首先，对非线性 Schr\"odinger 方程(\ref{eq:NLSE})的求解区域处理为离散的节点网格
\begin{align}
    u_{n}^{m}&=u(m\Delta\xi,n\Delta\tau)\\
    m&=0,1,2,\cdots,M \nonumber \\
    n&=-N,-N+1,\cdots,-1,0,1,\cdots,N-1,N \nonumber    
\end{align}
其中，上标 $m$ 标记传输距离，下标 $n$ 标记群速度坐标系下的时间。

其次，将非线性 Schr\"odinger 方程近似处理为等效的差分方程。一般应用 Taylor 级数对偏微分方程进行差分近似，比如 Euler 格式、蛙跳（hopscotch）格式、Crank-Nicolson 格式\cite{Taha,Qianshun}。1976 年，Ablowite 和 Ladik 基于逆散射方法提出了针对非线性偏微分方程求解的差分格式\cite{Taha}。

本文仅介绍常用的差分格式，最典型的差分近似是 Euler 差分格式
\begin{equation}
    i\frac{u_n^{m+1}-u_n^{m-1}}{2\Delta \xi}+\frac{u_{n+1}^{m}-2u_n^m+u_{n-1}^m}{2(\Delta\tau)^2}+|u_n^m|u_n^m=0
\end{equation}
Euler 格式是显式差分格式，一阶偏导 $\partial u/\partial \xi$ 为中心差分，二阶偏导 $\partial^2 u/\partial \tau^2$ 也是中心差分，因此其截断误差为 $O((\Delta\xi)^2)+O((\Delta\tau)^2)$。该差分格式的稳定性受 CFL 条件\footnote{该稳定性条件由 Courant、Friedrichs、Lewy 提出，可用 von Neumann 方法推导。}约束
\begin{equation}
    \lambda=\frac{\Delta\xi}{2(\Delta\tau)^2}\leqslant\frac{1}{4}
\end{equation}
当然也可以对非线性项做（时间）平均近似得到蛙跳格式
\begin{subequations}
    \begin{equation}
        i\frac{u_n^{m+1}-u_n^{m}}{\Delta\xi}+\frac{u_{n+1}^m-2u_n^m+u_{n-1}^m}{2(\Delta\tau)^2}+\frac{1}{2}\left(|u_{n+1}^m|^2u_{n+1}^m+|u_{n-1}^m|^2u_{n-1}^m\right)=0
        \label{eq:Hopscotch_explicit}
    \end{equation}
    \begin{equation}
        i\frac{u_n^{m+1}-u_n^{m}}{\Delta\xi}+\frac{u_{n+1}^{m+1}-2u_n^{m+1}+u_{n-1}^{m+1}}{2(\Delta\tau)^2}+\frac{1}{2}\left(|u_{n+1}^{m+1}|^2u_{n+1}^{m+1}+|u_{n-1}^{m+1}|^2u_{n-1}^{m+1}\right)=0
        \label{eq:Hopscotch_implicit}
    \end{equation}
\end{subequations}
把隐式差分格式(\ref{eq:Hopscotch_implicit})的上标 $m$ 代换为 $m-1$，并结合显式差分格式(\ref{eq:Hopscotch_explicit})可以发现蛙跳格式隐含了空间平均的近似
\begin{equation}
    u_n^{m+1}+u_n^{m-1}=2u_n^m
\end{equation}
因此，蛙跳格式的一阶偏导 $\partial u/\partial \xi$ 虽然是前向差分，但是其截断误差也是 $O((\Delta\xi)^2)+O((\Delta\tau)^2)$，并且无条件稳定。

比较常用的是一种差分近似是 Crank-Nicolson 格式，它将二阶偏导项 $\partial^2 u/\partial \tau^2$ 做了平均处理\cite{Greig,Taha}
\begin{equation}
    \begin{aligned}
        i\frac{u_n^{m+1}-u_n^m}{\Delta\xi}&+\frac{1}{4(\Delta\tau)^2}\left[u_{n+1}^{m+1}-2u_n^{m+1}+u_{n-1}^{m+1}+u_{n+1}^{m}-2u_n^{m}+u_{n-1}^{m}\right]\\
        &+\frac{1}{2}\left(|u_{n+1}^m|^2u_{n+1}^m+|u_{n-1}^m|^2u_{n-1}^m\right)=0
    \end{aligned}
\end{equation}
该种差分格式的截断误差为 $O(\Delta\xi)+O((\Delta\tau)^2)$，并且无条件稳定。若要减小 Crank-Nicolson 格式的空间截断误差，可对非线性项做（空间）平均
\begin{equation}
    \begin{aligned}
        i\frac{u_n^{m+1}-u_n^m}{\Delta\xi}&+\frac{1}{4(\Delta\tau)^2}\left[u_{n+1}^{m+1}-2u_n^{m+1}+u_{n-1}^{m+1}+u_{n+1}^{m}-2u_n^{m}+u_{n-1}^{m}\right]\\
        &+\frac{1}{2}\left(|u_n^{m+1}|^2u_n^{m+1}+|u_n^m|^2u_n^m\right)=0
    \end{aligned}
    \label{eq:Crank-Nicolson}
\end{equation}

为了更方便地在计算机上实现差分求解，通常把差分方程改写成矩阵方程的形式，比如对 Crank-Nicolson 格式的差分方程(\ref{eq:Crank-Nicolson})，可用矩阵方程表达\cite{Taha}
\begin{equation}
    \textbf{\textit{A}}\textbf{\textit{u}}^{m+1}=\textbf{\textit{F}}
    \label{eq:matrix}
\end{equation}
其中，差分矩阵 $\textbf{\textit{A}}$ 通常是一个带状矩阵，在周期性边界条件（$u_{-N-1}^m=u_N^m$，$u_{N+1}^m=u_{-N}^m$）的假设下
\begin{subequations}
    \begin{equation}
        \textbf{\textit{A}}=\begin{bmatrix}
            i-\lambda & \lambda/2 & & & \lambda/2\\
            \lambda/2 & i-\lambda &\lambda/2 & &\\
             & \ddots & \ddots & \ddots & \\
             & & \lambda/2 & i-\lambda & \lambda/2\\
             \lambda/2 & & & \lambda/2 & i-\lambda
        \end{bmatrix}\hspace{4ex}\textbf{\textit{u}}^{m+1}=\begin{bmatrix}
            u_{-N}\\
            u_{-N+1}\\
            \vdots\\
            u_{N-1}\\
            u_{N}
        \end{bmatrix}^{m+1}
    \end{equation}
    \begin{equation}
        \begin{aligned}
            F_n&=-\frac{\lambda}{2}\left(u_{n+1}^m+u_{n-1}^m\right)+(\lambda+i)u_n^m+\frac{\Delta\xi}{2}\left(|u_n^{m+1}|^2u_n^{m+1}+|u_n^m|^2u_n^m\right)\\
            &\approx-\frac{\lambda}{2}\left(u_{n+1}^m+u_{n-1}^m\right)+(\lambda+i)u_n^m+\Delta\xi|u_n^m|^2u_n^m                
        \end{aligned}
    \end{equation}
\end{subequations}

最后便是求解矩阵方程(\ref{eq:matrix})。差分矩阵 $\textbf{\textit{A}}$ 通常是大型矩阵，用 Gauss 消元法求解矩阵方程(\ref{eq:matrix})会花费较多的计算时间。一般应用迭代方法求解矩阵方程，常用的迭代算法有 Jacobi 迭代、Gauss-Seidel 迭代和逐次超松弛迭代（successive overrelaxation，简称 SOR）方法\cite{Anne}，本文不做具体介绍。在 MATLAB 中，应用 spdiags 函数可以很方便地创建差分矩阵 $\textbf{\textit{A}}$，并且 inv 函数可以直接求解带状矩阵 $\textbf{\textit{A}}$ 的逆矩阵，使得有限差分的编程更容易实现（详见附录3）。

有限差分法可以实现单光孤子周期性演化的计算，但是有限差分法耗时非常多，且不容易控制传输距离，很容易产生发散的结果（见图\ref{fig:FDM}）。一般情况下，分步 Fourier 方法比有限差分法快一到两个数量级，但是在波分复用系统的计算中，有限差分法比分步 Fourier 方法所用时间更少\cite{Agrawal}。
\begin{figure}[tbp]
    \centering
    \begin{subfigure}[t]{0.85\linewidth}
        \captionsetup{justification=centering} 
        \begin{minipage}[b]{1\linewidth}
        \includegraphics[width=0.45\linewidth]{FDM_基阶孤子_侧视图.png}
        \includegraphics[width=0.45\linewidth]{FDM_基阶孤子_俯视图.png}
        \caption{$u(0,\tau)=\mathrm{sech}(\tau)$}
        \end{minipage}
    \end{subfigure}\\
    \begin{subfigure}[t]{0.85\linewidth}
        \captionsetup{justification=centering} 
        \begin{minipage}[b]{1\linewidth}
        \includegraphics[width=0.45\linewidth]{FDM_二阶孤子_侧视图.png}
        \includegraphics[width=0.45\linewidth]{FDM_二阶孤子_俯视图.png}
        \caption{$u(0,\tau)=2\mathrm{sech}(\tau)$}
        \end{minipage}
    \end{subfigure}\\
    \begin{subfigure}[t]{0.85\linewidth}
        \captionsetup{justification=centering} 
        \begin{minipage}[b]{1\linewidth}
        \includegraphics[width=0.45\linewidth]{FDM_三阶孤子_侧视图.png}
        \includegraphics[width=0.45\linewidth]{FDM_三阶孤子_俯视图.png}
        \caption{$u(0,\tau)=3\mathrm{sech}(\tau)$}
        \end{minipage}
    \end{subfigure}\\
    \begin{subfigure}[t]{0.85\linewidth}
        \captionsetup{justification=centering} 
        \begin{minipage}[b]{1\linewidth}
        \includegraphics[width=0.45\linewidth]{FDM_四阶孤子_侧视图.png}
        \includegraphics[width=0.45\linewidth]{FDM_四阶孤子_俯视图.png}
        \caption{$u(0,\tau)=4\mathrm{sech}(\tau)$}
        \end{minipage}
    \end{subfigure}
    \caption{有限差分模拟的孤子周期性演化}
    \label{fig:FDM}
\end{figure}

\section{结论}
本文作者对伪谱方法和分步 Fourier 方法进行数值实验（有限差分法的计算时间比较长，笔者没有对其进行完整的数值实验），比较这两种数值方法的计算时间和求解结果（实验结果见附录4）
\begin{enumerate}[label=(\arabic*)]
    \item 从求解的结果看，伪谱方法的结果更加准确和稳定，分步 Fourier 方法需要更小的空间和时间步长才能达到满意的结果；
    \item 从计算耗时来看，分步 Fourier 方法比伪谱方法更快，主要原因是伪谱方法必须通过求解频域的常微分方程，再进行 Fourier 逆变换才能得到时域结果。
\end{enumerate}
对三种数值方法求解单光孤子问题进行总结（见表\ref{tab:compare}）
\begin{table}[htbp]
    \centering
    \caption{非线性 Schr\"odinger 方程的三种数值方法求解对比}
    \label{tab:compare}
    \begin{tabular}{cccc}
      \toprule 
        数值方法 & 计算类型 & 准确程度 & 计算时间 \\
      \midrule  
        伪谱方法 & 频域计算 & 准确 & 较快 \\
        分步 Fourier 方法 & 时域和频域交替计算 & 较准确 & 快 \\
        有限差分法 & 时域计算 & 较准确 & 非常慢 \\
      \bottomrule
    \end{tabular}
\end{table}

本文研究的问题是非线性 Schr\"odinger 方程的数值求解，着重研究单光孤子的周期性演化过程，在本文的最后对研究课题进行总结
\begin{enumerate}[label=(\arabic*)]
    \item 从解析方法看，非线性 Schr\"odinger 方程可以用 IST 方法求解；
    \item 从物理行为看，散射问题的离散本征值对应束缚态，其势场就是光孤子解；
    \item 从数值方法看，非线性 Schr\"odinger 方程可以用分步 Fourier 方法、伪谱方法、有限差分法进行求解。
\end{enumerate}
其实，非线性 Schr\"odinger 方程的研究方法、孤子概念都是从 KdV 方程迁移过来。就孤子问题来讲，还有很多有意思的课题，比如孤子稳定性、孤子相互作用等问题。




%%%%%%%%%%%%%%%%%%%%%%%%%%%%%%%%%%%%%%%%%%%%%%%%%%
% 临时标签，用于编译时追踪正文末尾
%%%%%%%%%%%%%%%%%%%%%%%%%%%%%%%%%%%%%%%%%%%%%%%%%%

%%%%%%%%%%%%%%%%%%%%%%%%%%%%%%%%%%%%%%%%%%%%%%%%%%
% 后续内容，标题三号黑体居中，章节无编号
% --------------------------------------------%

% https://www.zhihu.com/question/29413517/answer/44358389 %
% 说明如下：
% secnumdepth 这个计数器是 LaTeX 标准文档类用来控制章节编号深度的。在 article 中，这个计数器的值默认是 3，对应的章节命令是 \subsubsection。也就是说，默认情况下，article 将会对 \subsubsection 及其之上的所有章节标题进行编号，也就是 \part, \section, \subsection, \subsubsection。LaTeX 标准文档类中，最大的标题是 \part。它在 book 和 report 类中的层级是「-1」，在 article 类中的层级是「0」。这里，我们在调用 \appendix 的时候将计数器设置为 -2，因此所有的章节命令都不会编号了。不过，一般还是会保留 \part 的编号的。所以在实际使用中，将它设置为 0 就可以了。

% 在修改过程中请注意不要破环命令的完整性

\renewcommand\appendix{\setcounter{secnumdepth}{-2}}
\appendix

% 主文件有代码去掉页眉章节编号的“.”，但这会因为bug导致无编号章节显示一个错误编号，所以这里在无编号章节之前再次重定义sectionmark。
%\renewcommand{\sectionmark}[1]{\markright{#1}}

% section 标题从这里往后改为三号居中
\titleformat{\section}{\centering \bfseries \zihao{3}}{\thesection}{1em}{}


%!TEX root = ../GLUTthesis.tex
\section{致\quad 谢} % 无章节编号

% 公元二零二零年，岁次庚子，仲夏之月，吾之论文乃告杀青。辞穷理微，未敢称凌云之作，镂心鸟迹，得不效相如之叹?于是凭窗抱膝，寄情遐思。忆吾弱冠之龄入理工大学，意气方遒。尔来春秋有四，于今毕业，年齿已趋而立。户牅之外，万物滋荣，景致阙如昨日，堂室之内，吾已有苍颜白发矣。文凭两纸霜鬓两行，黄粱一枕功名一场，此皆寻常人生，乏善可陈。然吾身发受之父母，道德受之母校，学问受之师长，育教之恩，虽陨首结草不能报之万一。是以为情造文，铭而致谢。

大学四年即将结束，我也快要完成我的毕业论文，首先要感谢支持我完成论文的朋友。感谢我的指导老师王恒教授，给了我很大的发挥空间，在我对课题毫无头绪的时候指导我确定了具体的研究方向；感谢提供论文 \LaTeX 模板的胡光辉老师，他设计的论文模板让我不再被公式输入、排版等问题所困扰；感谢提供本文研究课题的巴灵丽同学，光孤子数值模拟问题来自她参加的夏令营的一个题目；感谢提供计算算力的裴俊杰同学，没有他的算力支持，我得不到有限差分的计算结果$\ldots\ldots$

其次要感谢生我养我的父母。我能来到这个世界这本身就是一件幸运的事，更何况是他们无微不至的关爱让我可以顺利健康地长大成人。小时候很淘气，给父母惹了很多麻烦；青春期比较叛逆，觉得父母老古董了，经常和他们对着干；长大了才发现，维持一个家庭，上有老，下有小是多么的不容易。父母正在老去，我要更加努力，为自己和社会创造更多价值，让他们不再操劳，能安度晚年。

最后要感谢陪伴我四年的桂林理工大学。感谢理学院的老师和同学，他们让我大学四年的学习生活不再孤单，经常给我提供帮助和建议以及挑战；感谢田径队的教练和队员，在大学因为共同的跑步爱好而聚在一起，互相鼓励、交流、提高，真的非常难得；感谢电子协会、数学建模协会和物理协会，让我大学四年的课余生活丰富多彩，社团活动给了我很多挑战，让我练出了非常多的技能$\ldots\ldots$即将离开母校，回想这四年，虽然有很多难熬的、烦心的事情，但是印象最深的还是那些奋斗过，快乐过的日子。
\newpage


% \section{参考文献} % bibliography会自动显示参考文献四个字
\addcontentsline{toc}{section}{参考文献} % 由于参考文献不是section，这句把参考文献加入目录
% \nocite{*} % 该命令用于显示全部参考文献，即使文中没引用
% cls文件中已经引入package，这里不需要调用 \bibliographystyle 了。
% \bibliographystyle{gbt7714-2005}
% \bibliography{thesis-references}
\printbibliography
\newpage

% 致谢应该在参考文献的后面，不编章号，而附录应该位于致谢后。


%!TEX root = ../GLUTthesis.tex
\section{附录1} % 无章节编号
在 MATLAB 上实现用伪谱方法求解非线性 Schr\"odinger 方程
\begin{lstlisting}[language=Matlab,
  morekeywords={},
  emph={fft, ifft, ode45},
  caption={伪谱方法求解三阶孤子}
  ]
    %% 参数设置
    sgn = -1;
    T_max = 10; N_t = 256; 
    t = 2*T_max/N_t*[-N_t/2:N_t/2-1];
    omega = (pi/T_max)*[0:N_t/2-1 -N_t/2:-1].';
    Z = 4; h = 0.2;
    z = 0:h:Z;

    %% 初始条件
    u = 3*sech(t); ut = fft(u);

    %% ODE求解
    [z, UT] = ode45('NLSE', z, ut, [], sgn, omega);
    U = ifft(UT, [], 2);

    %% 可视化
    figure
    waterfall(t, z, abs(U));
    xlabel('时间T/T_0'), ylabel('传输距离z/L_D'), zlabel('强度')
    figure
    waterfall(fftshift(omega)./(2.*pi), z, abs(fftshift(UT, 2)));
    xlabel('频率(\nu-\nu_0)T_0'), ylabel('传输距离z/L_D'), zlabel('频谱强度')
\end{lstlisting}
其中，调用 ode45 时需要自定义表达常微分方程的函数
\begin{lstlisting}[language=Matlab,
  morekeywords={},
  emph={fft, ifft},
  caption={}
  ]
    function dut = NLSE(z, ut, dummy, sgn, omega)
    %NLSE 函数表示被求解的 ODE 方程
    %   dut = i*sgn/2*(omega^2)*ut+i*F[|u|^2*u]
    u = ifft(ut);
    dut = sgn*(i/2)*(omega.^2).*ut+i*fft((abs(u).^2).*u);
    end
\end{lstlisting}

\newpage
\section{附录2}
在 MATLAB 上实现用分步 Fourier 方法求解非线性 Schr\"odinger 方程
\begin{lstlisting}[language=Matlab,
  morekeywords={},
  emph={fft, ifft, ode45},
  caption={分步 Fourier 方法求解三阶孤子}
  ]
    %% 参数设置
    beta_2 = -1; N = 1;
    T_max = 10; N_t = 512;
    Z = 4; h = 0.01;
    t = 2*T_max/N_t*[-N_t/2:N_t/2-1];
    omega = (pi/T_max).*[0:N_t/2-1 -N_t/2:-1];
    z = 0:h:Z;
    
    %% 色散项和非线性项
    D = exp(1i/2*beta_2*omega.^2*h);
    NL = 1i*N^2*h;
    
    %% 初始条件
    u = 3*sech(t);
    U = []; UT = [];
    U(1,:) = u; UT(1,:) = fftshift(fft(u));
    
    %% 分步 Fourier 求解 1/2N->D->1/2N
    u = u.*exp(abs(u).^2.*NL/2);
    for n = 2:length(z)
        ut = fft(u).*D;
        u = ifft(ut);
        u = u.*exp(abs(u).^2.*NL);
        u_r = u.*exp(-abs(u).^2.*NL/2);
        U(n,:) = u_r;
        UT(n,:) = fftshift(fft((u_r)));
    end

    %% 可视化
    figure
    waterfall(t, z, abs(U));
    xlabel('时间T/T_0'), ylabel('传输距离z/L_D'), zlabel('强度')
    figure
    waterfall(fftshift(omega)./(2.*pi), z, abs(UT));
    xlabel('频率(\nu-\nu_0)T_0'), ylabel('传输距离z/L_D'), zlabel('频谱强度')    
\end{lstlisting}

\newpage
%%%%%%%%%%%%%%%%%%%%%%%%%%%%%%%%%
\section{附录3}
在 MATLAB 上实现用有限差分法求解非线性 Schr\"odinger 方程
\begin{lstlisting}[language=Matlab,
  morekeywords={},
  emph={spdiags,inv},
  caption={有限差分法求解三阶孤子}
  ]
  %% 参数设置
  Z_max = 4; T_max = 10;
  dz = 0.0001; dt = 0.002; lambda = dz/dt^2/2;
  z = [0:dz:Z_max]; t = [-T_max:dt:T_max];
  [T, Z] = meshgrid(t, z);
  N_t = length(t); N_z = length(z);
  
  %% 差分矩阵
  e = ones(N_t, 1);
  A = spdiags([lambda/2*e, (1i-lambda)*e, lambda/2*e], [-1, 0, 1], N_t, N_t);
  A(1, N_t) = lambda/2; A(N_t, 1) = lambda/2;
  B = spdiags([-lambda/2*e, (1i+lambda)*e, -lambda/2*e], [-1, 0, 1], N_t, N_t);
  B(1, N_t) = -lambda/2; B(N_t, 1) = -lambda/2;
  A_i = inv(A);
  
  %% 边界条件
  u = 3*sech(t)'; U(1,:) = u';
  
  %% 差分求解
  for m = 2:N_z
      F = B*u-dz*abs(u).^2.*u;
      u = A_i*F;
      U(m,:) = u';
  end
  
  %% 可视化
  figure
  surf(T, Z, abs(U));
  xlabel('时间T/T_0'), ylabel('传输距离z/L_D'), zlabel('强度')
\end{lstlisting}

\newpage
%%%%%%%%%%%%%%%%%%%%%%%%%%%%%%%%%
\section{附录4}
\begin{table}[htbp]
  \centering
  \caption*{伪谱方法的数值实验结果（耗时记录）}
  \begin{tabular}{lcccccc}
    \toprule 
      $\Delta \xi$ & $N_\tau$ & 基阶孤子/s & 二阶孤子/s & 三阶孤子/s & 四阶孤子/s & 五阶孤子/s \\
    \midrule  
      0.1 & 256 & 0.913524 & 0.867968 & 1.087755 & 结果发散 & 结果发散 \\
      0.05 & 256 & 0.8847945 & 0.904476 & 1.126649 & 结果发散 & 结果发散 \\
      0.03 & 256 & 0.906814 & 0.887762 & 1.105801 & 结果发散 & 结果发散 \\
      0,01 & 256 & 0.961330 & 0.959687 & 1.207335 & 结果发散 & 结果发散 \\
      0.005 & 256 & 1.032782 & 1.017613 & 1.255176 & 结果发散 & 结果发散 \\
      0.003 & 256 & 1.128874 & 1.152278 & 1.403453 & 结果发散 & 结果发散 \\
      0.001 & 256 & 1.605832 & 1.594857 & 1.800561 & 结果发散 & 结果发散 \\
      0.1 & 512 & 4.150410 & 4.064225 & 4.060273 & 4.303629 & 5.136246 \\
      0.05 & 512 & 4.163852 & 4.063956 & 4.088162 & 4.283738 & 5.169811 \\
      0.03 & 512 & 4.083730 & 4.160560 & 4.200042 & 4.375078 & 5.205791 \\
      0.01 & 512 & 4.245213 & 4.231423 & 4.255282 & 4.445187 & 5.340950 \\
      0.005 & 512 & 4.390212 & 4.398500 & 4.402081 & 4.643569 & 5.505080 \\
      0.003 & 512 & 4.787179 & 4.706263 & 4.750699 & 4.907478 & 5.813890 \\
      0.001 & 512 & 5.457064 & 5.476309 & 5.466284 & 5.646329 & 6.594705 \\
    \bottomrule
    \end{tabular}
\end{table}

\begin{table}[tbp]
  \centering
  \caption*{分步 Fourier 方法的数值实验结果（耗时记录）}
  \begin{tabular}{lcccccc}
    \toprule 
      $\Delta \xi$ & $N_\tau$ & 基阶孤子/s & 二阶孤子/s & 三阶孤子/s & 四阶孤子/s & 五阶孤子/s \\
    \midrule
      0.1 & 256 & 0.011991 & 0.011326 & 结果发散 & 结果发散 & 结果发散 \\
      0.05 & 256 & 0.012258 & 0.012525 & 结果发散 & 结果发散 & 结果发散 \\
      0.03 & 256 & 0.019958 & 0.015117 & 结果发散 & 结果发散 & 结果发散 \\
      0.01 & 256 & 0.047944 & 0.048051 & 0.040772 & 0.046679 & 结果发散 \\
      0.005 & 256 & 0.098894 & 0.101310 & 0.102865 & 0.101067 & 结果发散 \\
      0.003 & 256 & 0.438257 & 0.431412 & 0.416631 & 结果发散 & 结果发散 \\
      0.001 & 256 & 4.911769 & 4.917287 & 4.916656 & 4.927861 & 结果发散 \\
      0.1 & 512 & 0.010518 & 0.010189 & 结果发散 & 结果发散 & 结果发散 \\
      0.05 & 512 & 0.013548 & 0.015026 & 结果发散 & 结果发散 & 结果发散 \\
      0.03 & 512 & 0.020160 & 0.019286 & 结果发散 & 结果发散 & 结果发散 \\
      0.01 & 512 & 0.066089 & 0.059894 & 0.058442 & 结果发散 & 结果发散 \\
      0.005 & 512 & 0.362146 & 0.371222 & 0.359213 & 结果发散 & 结果发散 \\
      0.003 & 512 & 1.137290 & 1.163925 & 1.134111 & 1.150079 & 1.142337 \\
      0.001 & 512 & 9.897807 & 9.908339 & 9.984559 & 9.904616 & 9.893836 \\
    \bottomrule
    \end{tabular}
\end{table}

\newpage
\section*{MATLAB 编程要点}
\subsection*{ODE 求解器的调用}
在 MATLAB 中调用 ODE 求解器通常需要一个调用指令 odesolver 和一个函数定义文件 odefun.m。{\bfseries odesolver 的调用语法为}
$$
\mathrm{[X,U]=odesolver(odefun, xspan, u0, options, parameter1, parameter2, \ldots)}
$$
\begin{enumerate}[label=(\arabic*)]
  \item “odesolver”是所调用求解器的函数名称，可以是 ode45、ode23、ode113、ode15s、ode23s、ode23t、ode23tb 和 ode15i 中的一个。
  \item “odefun”是定义常微分方程组 $u'_i=f_i(x,u_1,u_2,\ldots,u_n)$ 的函数句柄，可以是 odefun.m 文件定义的函数，也可以是“@”标记的匿名函数。
  \item “xspan”是表示常微分方程组自变量 $x$ 取值范围的向量 $[x_0,x_1,x_2,\ldots,x_n]$。
  \item “u0”是代表常微分方程组初始条件的向量 $u_i(x_0)$。
  \item “options”是一个改变 odesolver 默认参数设置的结构体，是一个可选参数。若不改变默认参数设置，可将其设置为空。
  \item “parameter1, parameter2,\ldots”是传递给 odefun 的参数，是可选项。
  \item “[X,U]”保存计算结果。向量 X 就是自变量的离散取值 $[x_0,x_1,x_2,\ldots,x_n]$，矩阵 U 就是所求常微分方程组的数值解。
\end{enumerate}

{\bfseries odefun 函数的定义需要与 odesolver 的调用形式相对应}
$$
\mathrm{function\; du\; =\; odefun(x, u, dummy,  parameter1, parameter2, \ldots)}
$$
\begin{enumerate}[label=(\arabic*)]
  \item “odefun”为函数名，函数文件名必须与其相同。
  \item “x”为自变量。
  \item “u”是常微分方程组的函数值 $u_i(x)$，是一个矩阵。
  \item “dummy”是一个占位符，无实际意义。
  \item “parameter1, parameter2, \ldots”用于接收传递来的参数。
  \item “du”为函数输出结果的变量名。
\end{enumerate}

\subsection*{fftshift 操作}
由于 FFT 算法本身的原因，fft 函数的输出结果在频域上是颠倒的
$$
0\;,2\pi/T_0\;,4\pi/T_0,\ldots,(2N-4)\pi/T_0\;,(2N-2)\pi/T_0\;,-2N\pi/T_0\;,-(2N-2)\pi/T_0,\ldots,-4\pi/T_0\;,-2\pi/T_0
$$
因此需要 fftshift 函数调整 fft 输出序列的顺序。但是在分步 Fourier 方法（伪谱方法）求解问题时，不需要每次运算都调整频谱顺序，只需将最后一步运算的频谱结果调整即可（频率序列首先需要人为颠倒，可视化时再调整回来）。

\subsection*{spdiags 函数创建带状矩阵}
$$
\mathrm{A\; =\; spdiags(B,d,m,n)}
$$
通过获取 B 的列并沿 d 指定的对角线放置它们，来创建一个 m×n 带状矩阵。

\subsection*{可视化模板}
\begin{multicols}{2}
  \begin{lstlisting}[language=Matlab,
  morekeywords={},
  emph={waterfall, view, colormap},
  caption={}
  ]
    %% 瀑布图可视化
    figure
    waterfall(t, z, abs(U))
    view(-45,45)
    xlabel('时间 T/T_0'),
    ylabel('传输距离 z/L_D'),
    zlabel('强度'),
    axis([-10 10 0 inf 0 inf]),
    colormap([0 0 0]); 
  \end{lstlisting}
  \begin{figure}[H]
    \centering
    \includegraphics[width=0.35\textwidth]{瀑布图可视化.png}
    % \caption{瀑布图可视化}
  \end{figure}
\end{multicols}

\begin{multicols}{2}
  \begin{lstlisting}[language=Matlab,
  morekeywords={},
  emph={meshgrid, surf, shading, view, colormap},
  caption={}
  ]
    %% 曲面图可视化
    [T, Z] = meshgrid(t, z);
    figure
    surf(T, Z, abs(U))
    xlabel('时间 T/T_0'),
    ylabel('传输距离 z/L_D'),
    zlabel('强度'),
    view(-20,70)
    axis([-10 10 0 inf 0 inf]),
    shading interp
    colormap(gca, 'jet');
  \end{lstlisting}
  \begin{figure}[H]
    \centering
    \includegraphics[width=0.35\textwidth]{曲面图可视化.png}
    % \caption{曲面图可视化}
  \end{figure}
\end{multicols}


\end{document}
